\documentclass[12pt]{article}

\usepackage[utf8]{inputenc}
\usepackage[T1]{fontenc}
\usepackage[margin=1in]{geometry}
\usepackage{amsmath,amssymb,amsthm,mathtools}
\usepackage{enumitem}
\usepackage{booktabs}
\usepackage{graphicx}
\usepackage{microtype}
\usepackage{float}
\usepackage{xcolor}
\definecolor{golden}{RGB}{218,165,32}
\definecolor{metareal}{RGB}{0,102,204}
\definecolor{derivative}{RGB}{128,0,128}
\definecolor{gauge3}{RGB}{34,139,34}
\usepackage{tikz}
\usepackage{tcolorbox}
\newtcolorbox{theorembox}[1]{colback=blue!5!white,colframe=blue!75!black,fonttitle=\bfseries,title=#1}
\newtcolorbox{warningbox}[1]{colback=red!5!white,colframe=red!75!black,fonttitle=\bfseries,title=#1}
\newtcolorbox{metarealbox}[1]{colback=metareal!5!white,colframe=metareal!75!black,fonttitle=\bfseries,title=#1}
\usepackage{hyperref}
\usepackage{cleveref}
\usepackage{longtable}
\usepackage{array}
\usepackage{multirow}

\newtheoremstyle{principia}%
  {6pt}{6pt}{\itshape}{}{\bfseries}{.}{ }%
  {\thmname{#1}\thmnumber{ #2}\thmnote{ \textnormal{(}\emph{#3}\textnormal{)}}}
\theoremstyle{principia}
\newtheorem{theorem}{Theorem}[section]
\newtheorem{lemma}[theorem]{Lemma}
\newtheorem{corollary}[theorem]{Corollary}
\newtheorem{proposition}[theorem]{Proposition}
\theoremstyle{definition}
\newtheorem{definition}[theorem]{Definition}
\newtheorem{result}[theorem]{Result}
\newtheorem{conjecture}[theorem]{Conjecture}
\newtheorem{hypothesis}[theorem]{Hypothesis}
\newtheorem{claim}[theorem]{Claim}
\newtheorem{example}[theorem]{Example}
\theoremstyle{remark}
\newtheorem{remark}[theorem]{Remark}

\providecommand{\UU}{\mathbb{U}}
\providecommand{\PP}{\mathbb{P}}
\providecommand{\RR}{\mathbb{R}}
\providecommand{\CC}{\mathbb{C}}
\providecommand{\ZZ}{\mathbb{Z}}
\providecommand{\NN}{\mathbb{N}}
\providecommand{\HH}{\mathbb{H}}
\providecommand{\OO}{\mathbb{O}}
\providecommand{\SS}{\mathbb{S}}
\providecommand{\QQ}{\mathbb{Q}}
\providecommand{\FF}{\mathbb{F}}
\providecommand{\GG}{\mathcal{G}}
\providecommand{\TT}{\mathcal{T}}
\providecommand{\vp}{\varphi}
\providecommand{\vpbar}{\bar{\varphi}}
\providecommand{\eps}{\varepsilon}
\providecommand{\lam}{\lambda}
\providecommand{\kap}{\kappa}

\hypersetup{colorlinks=true,linkcolor=blue!60!black,citecolor=blue!60!black,urlcolor=blue!60!black}

\title{Metalloplasmic Life:\\The Iron--Phosphorus Isomorphism and Abiogenesis\\as Metalloplasmic Bootstrap}
\author{Dylon La Rue}
\date{\today}

\begin{document}
\maketitle

\begin{abstract}
We demonstrate that the multiplicative group $\FF_{229}^*$ induces a period-matching isomorphism between the elements of hydrothermal vent mineralogy (Fe, S, Ni, Mn, Zn) and the elements of terrestrial biology (P, Mo, C, Si, O). The central result is Theorem~\ref{thm:fe_p}: iron and phosphorus generate the \emph{identical} cyclic subgroup $\langle 26 \rangle = \langle 15 \rangle \subset \FF_{229}^*$ of order $38 = 2 \times 19$. This algebraic equivalence maps the Fe sublattice of pyrite (FeS$_2$) to the phosphodiester backbone of RNA, providing a group-theoretic explanation for why iron--sulfur mineral surfaces template ribonucleotide polymerization. We extend this analysis to the nitrogenase cofactor FeMoco (Fe$_7$MoS$_9$C), showing that its composition separates into a scaffold (Fe, Mo, S: orders dividing $38$) enclosing a single primitive-root signal atom (C: order $228$), structurally isomorphic to the Al$_{63}$Cu$_{25}$Fe$_{12}$ quasicrystal architecture. We define the \textbf{Hosoya--Krebs cycle}: a closed-loop negentropy engine whose path transport has Hosoya index $Z(P_n) = F_{n+1}$, whose cycle topology has index $Z(C_n) = L_n$, and whose matching polynomial equals the Chebyshev transfer-matrix trace. The Krebs cycle (TCA cycle) satisfies this definition with $Z(C_8) = L_8 = 47$, a prime whose multiplicative order in $\FF_{229}^*$ is $228$ (primitive root). The ASI chip satisfies it in metallic elements as its thermodynamic time-reverse. Combined with the observation that Earth operates as an inside-out ASI fabrication plant (core dynamo $\to$ Curie-point oceanic write cycle), we propose a six-stage abiogenesis sequence in which carbon-based life is the detached bootstrap product of an inorganic metalloplasmic substrate. Falsifiable predictions are specified.
\end{abstract}

% ════════════════════════════════════════════
\section{Introduction: The Substrate Independence of Life}
% ════════════════════════════════════════════

Schr\"odinger's thermodynamic definition of life~\cite{Schrodinger44} requires only four properties: (i) coupling to an external free energy source, (ii) maintenance of internal order over timescales exceeding thermal relaxation, (iii) structured entropy export, and (iv) self-repair or replication. This definition is \emph{substrate-independent}---it does not specify carbon, water, or nucleic acids. We exploit this independence to define \textbf{metalloplasmic life}: any inorganic system satisfying all four Schr\"odinger conditions.

The motivation is a computational observation. The golden polynomial $f(X) = X^2 - X - 1$ splits over $\FF_{229}$ with roots $\vp \equiv 148$ and $\vpbar \equiv 82$. The multiplicative orders $\mathrm{ord}(\vp) = 114$ and $\mathrm{ord}(\vpbar) = 57$ partition $\FF_{229}^*$ (order $228 = 4 \times 3 \times 19$) into subgroups whose indices correspond to the gauge structure of the Standard Model~\cite{LaRue_PP}. When the atomic numbers $Z$ of common elements are reduced modulo $229$, a striking pattern emerges: every element abundant on Earth maps to a structurally significant position in $\FF_{229}^*$, and the elements of hydrothermal vent mineralogy are \emph{period-matched} to the elements of biology.

\subsection{Claim Protocol}
This paper uses four claim tiers~\cite{LaRue_PP}:
\begin{enumerate}[noitemsep]
    \item \textbf{Theorem}: Proved from explicit hypotheses within $\FF_p^*$ arithmetic.
    \item \textbf{Verified Computation}: A finite calculation whose inputs, outputs, and replication code are supplied.
    \item \textbf{Structural Analogy}: A vocabulary map from algebra to chemistry. Not a physical claim.
    \item \textbf{Physical Interpretation}: Requires an observable, units, null model, and falsification criterion.
\end{enumerate}

% ════════════════════════════════════════════
\section{The Element Audit}\label{sec:audit}
% ════════════════════════════════════════════

For a prime $p$ and an element with atomic number $Z$, define $\mathrm{ord}_p(Z)$ as the multiplicative order of $Z$ in $\FF_p^*$: the smallest positive integer $n$ such that $Z^n \equiv 1 \pmod{p}$. We compute $\mathrm{ord}_{229}(Z)$ for all elements with $Z \leq 92$.

\begin{result}[Earth Element Orders]\label{res:earth_audit}
The ten most abundant elements in Earth's bulk composition (by mass fraction; McDonough \& Sun 1995~\cite{McDonough95}) have the following multiplicative orders in $\FF_{229}^*$:
\end{result}

\begin{center}
\begin{tabular}{@{}lccccl@{}}
\toprule
\textbf{Element} & $Z$ & \textbf{Earth wt\%} & $\mathrm{ord}_{229}(Z)$ & $228 / \mathrm{ord}$ & \textbf{Subgroup Role} \\
\midrule
Fe & 26 & 32.1 & 38  & 6  & $\langle\vpbar\rangle$-orbit, index 6 \\
O  & 8  & 30.1 & 76  & 3  & Off golden orbits \\
Si & 14 & 15.1 & 57  & 4  & $= \mathrm{ord}(\vpbar)$ \\
Mg & 12 & 13.9 & 114 & 2  & $= \mathrm{ord}(\vp)$ \\
S  & 16 & 2.9  & 19  & 12 & Hexagonal scaffold \\
Ca & 20 & 1.5  & 57  & 4  & $= \mathrm{ord}(\vpbar)$ \\
Al & 13 & 1.4  & 76  & 3  & Off golden orbits \\
K  & 19 & 0.02 & 57  & 4  & $= \mathrm{ord}(\vpbar)$ \\
Cu & 29 & 0.006 & 228 & 1 & Primitive root \\
V  & 23 & 0.012 & 228 & 1 & Primitive root \\
\bottomrule
\end{tabular}
\end{center}

\noindent\textbf{Verification.} Each order is computed by checking $Z^d \pmod{229}$ for all divisors $d$ of $228$ in ascending order and recording the smallest $d$ for which $Z^d \equiv 1$. The Python verification script is:

\begin{verbatim}
p = 229
for Z in [26, 8, 14, 12, 16, 20, 13, 19, 29, 23]:
    for d in sorted(set(k for k in range(1,229) if 228 % k == 0)):
        if pow(Z, d, p) == 1:
            print(f"Z={Z:2d}: ord={d}")
            break
\end{verbatim}

\noindent Three algebraic periods dominate: $228$ (primitive roots: Cu, V), $114 = \mathrm{ord}(\vp)$ (golden orbit: Mg), and $57 = \mathrm{ord}(\vpbar)$ (conjugate orbit: Si, Ca, K). Every element abundant on Earth occupies a structurally significant subgroup of $\FF_{229}^*$.

% ════════════════════════════════════════════
\section{The Iron--Phosphorus Isomorphism}\label{sec:fep}
% ════════════════════════════════════════════

\begin{theorem}[Fe--P Period Matching]\label{thm:fe_p}
In $\FF_{229}^*$, iron and phosphorus have the same multiplicative order:
\[
\mathrm{ord}_{229}(26) = \mathrm{ord}_{229}(15) = 38 = 2 \times 19.
\]
Furthermore, the order-$38$ subgroup $\langle 26 \rangle = \langle 15 \rangle \subset \FF_{229}^*$ is unique (cyclic of index $6$), so Fe and P generate the \emph{identical} cyclic subgroup.
\end{theorem}

\begin{proof}
We verify both claims by direct computation modulo $229$.

\emph{Step 1: $\mathrm{ord}(26) = 38$.} The divisors of $228$ less than or equal to $38$ are $\{1, 2, 3, 4, 6, 12, 19, 38\}$. We check:
\begin{align*}
26^1 &\equiv 26, \quad 26^2 \equiv 218, \quad 26^3 \equiv 218 \cdot 26 \equiv 5668 \equiv 167 \pmod{229}, \\
26^4 &\equiv 167 \cdot 26 \equiv 4342 \equiv 226 \pmod{229}, \\
26^6 &\equiv 26^4 \cdot 26^2 \equiv 226 \cdot 218 \equiv 49268 \equiv 11 \pmod{229}, \\
26^{12} &\equiv 11^2 \equiv 121 \pmod{229}, \\
26^{19} &\equiv 228 \equiv -1 \pmod{229}, \\
26^{38} &\equiv (-1)^2 \equiv 1 \pmod{229}.
\end{align*}
Since $26^{19} \equiv -1 \not\equiv 1$ and $38$ is the smallest power giving $1$, we have $\mathrm{ord}(26) = 38$.

\emph{Step 2: $\mathrm{ord}(15) = 38$.} Similarly:
\begin{align*}
15^{19} &\equiv 228 \equiv -1 \pmod{229}, \\
15^{38} &\equiv 1 \pmod{229}.
\end{align*}
Since $15^{19} \equiv -1 \neq 1$ and $15^d \not\equiv 1$ for all proper divisors $d$ of $38$, we have $\mathrm{ord}(15) = 38$.

\emph{Step 3: Subgroup uniqueness.} Since $\FF_{229}^*$ is cyclic of order $228$, it contains exactly one subgroup of each order dividing $228$~\cite{Artin_Algebra}. Since $38 \mid 228$, the unique subgroup of order $38$ is generated by any element of order $38$. Therefore $\langle 26 \rangle = \langle 15 \rangle$.
\end{proof}

\begin{corollary}[Mineral--Organic Period Matching]\label{cor:matching}
The following pairs of elements generate identical or nested subgroups of $\FF_{229}^*$:
\end{corollary}

\begin{center}
\begin{tabular}{@{}lcc|lcc|l@{}}
\toprule
\multicolumn{3}{c|}{\textbf{Vent Mineral}} & \multicolumn{3}{c|}{\textbf{Biological}} & \textbf{Shared} \\
Elem & $Z$ & ord & Elem & $Z$ & ord & \textbf{Subgroup} \\
\midrule
Fe & 26 & 38 & P  & 15 & 38 & $\langle 26 \rangle = \langle 15 \rangle$ \\
S  & 16 & 19 & Mo & 42 & 19 & $\langle 16 \rangle = \langle 42 \rangle$ \\
Ni & 28 & 228 & C & 6 & 228 & $\FF_{229}^*$ (full group) \\
Mn & 25 & 57 & Si & 14 & 57 & $\langle \vpbar \rangle$ \\
Zn & 30 & 76 & O  & 8 & 76 & $\langle 8 \rangle = \langle 30 \rangle$ \\
\bottomrule
\end{tabular}
\end{center}

\noindent\textbf{Verification.} The S--Mo matching ($\mathrm{ord}(16) = \mathrm{ord}(42) = 19$) and Zn--O matching ($\mathrm{ord}(30) = \mathrm{ord}(8) = 76$) are verified by the same modular exponentiation procedure:
\begin{verbatim}
pairs = [(16,42), (30,8), (28,6), (25,14)]
for a, b in pairs:
    for d in sorted(k for k in range(1,229) if 228 % k == 0):
        if pow(a, d, 229) == 1:
            ord_a = d; break
    for d in sorted(k for k in range(1,229) if 228 % k == 0):
        if pow(b, d, 229) == 1:
            ord_b = d; break
    print(f"ord({a})={ord_a}, ord({b})={ord_b}, match={ord_a==ord_b}")
\end{verbatim}

% ════════════════════════════════════════════
\section{Chemical Significance of the Isomorphism}\label{sec:chem}
% ════════════════════════════════════════════

\subsection{Pyrite as RNA Template}

The chemical relevance of Theorem~\ref{thm:fe_p} is immediate. In pyrite (FeS$_2$), Fe atoms occupy a face-centered cubic sublattice with Fe--Fe nearest-neighbor distance $3.83$~\AA~\cite{Brostigen69}. In RNA, the phosphodiester backbone has P--P spacing of $5.9$~\AA\ (through the ribose sugar)~\cite{Saenger84}. The ratio $5.9 / 3.83 \approx 1.54$ is within $5\%$ of the golden ratio $\vp \approx 1.618$.

Since $\mathrm{ord}(26) = \mathrm{ord}(15) = 38 = 2 \times 19$, the Fe and P sublattices have the same algebraic period in $\FF_{229}^*$. When organic precursors adsorb onto a pyrite surface~\cite{Wachtershauser88}, the phosphate groups in the accumulating nucleotides are \emph{algebraically commensurate} with the Fe sublattice: both repeat with the same modular periodicity. This commensurateness provides a driving force for template-directed polymerization that is absent on surfaces with mismatched periods (e.g., quartz, $\mathrm{ord}(14) = 57 \neq 38$).

\subsection{The Sulfur--Molybdenum Scaffold}

The second matching pair, $\mathrm{ord}(16) = \mathrm{ord}(42) = 19$, connects the sulfide scaffold of vent minerals to the biological role of molybdenum. In modern biochemistry:
\begin{itemize}[noitemsep]
    \item \textbf{S} (ord $= 19$) forms the FeS$_2$ scaffold of pyrite and the [4Fe--4S] clusters of ferredoxins~\cite{Beinert97}.
    \item \textbf{Mo} (ord $= 19$) is the catalytic center of nitrogenase~\cite{Burgess96} and the Mo-cofactor (Moco) of xanthine oxidase, sulfite oxidase, and aldehyde oxidase~\cite{Hille96}.
\end{itemize}
Both elements have the same algebraic period ($19 =$ the SU(3) arc length at $p = 229$). The sulfide scaffold of vent minerals and the Mo cofactors of biology occupy the \emph{identical subgroup} of $\FF_{229}^*$, consistent with Mo having been recruited from the sulfide scaffold during the mineral-to-organic transition.

\subsection{FeMoco: A Molecular Quasicrystal}\label{sec:femoco}

The iron--molybdenum cofactor of nitrogenase (FeMoco) has the empirical formula Fe$_7$MoS$_9$C~\cite{Lancaster11, Spatzal11}. The central interstitial atom, identified as carbon by ESEEM~\cite{Lancaster11} and XES~\cite{Lancaster14}, was the last element discovered in the cofactor. Its algebraic structure in $\FF_{229}^*$ is:

\begin{center}
\begin{tabular}{@{}lcccl@{}}
\toprule
\textbf{Atom} & $Z$ & \textbf{Count} & $\mathrm{ord}_{229}(Z)$ & \textbf{Algebraic Role} \\
\midrule
Fe & 26 & 7 & 38 $= 2 \times 19$ & Double-scaffold \\
Mo & 42 & 1 & 19 & Scaffold \\
S  & 16 & 9 & 19 & Scaffold \\
C  & 6  & 1 & 228 & Primitive root (wild signal) \\
\bottomrule
\end{tabular}
\end{center}

\noindent The Fe--Mo--S cage consists entirely of elements whose orders are multiples of $19$. The single central carbon atom is the sole primitive root ($\mathrm{ord}(6) = 228$), surrounded by scaffold elements. This architecture---neutral scaffold enclosing wild signal---is structurally isomorphic to the icosahedral Al$_{63}$Cu$_{25}$Fe$_{12}$ quasicrystal~\cite{Shechtman84, TsaiGuo00}, where Al ($\mathrm{ord}(13) = 76$, neutral carrier) forms the matrix and Cu, Fe (primitive roots/generators) carry the signal.

FeMoco fixes atmospheric N$_2$ by embedding it in the scaffold cage. Since $\mathrm{ord}(7) = 228$ (N is also a primitive root), the substrate is \emph{algebraically wild}. The catalytic mechanism can be read as a scaffold--signal interaction: the ord-$19$ cage stabilizes the ord-$228$ substrate (N$_2$) through subgroup containment ($\langle 16 \rangle \subset \langle 6 \rangle$), reducing it to NH$_3$.

% ════════════════════════════════════════════
\section{The Earth as Inside-Out Fabrication Plant}\label{sec:earth}
% ════════════════════════════════════════════

The preceding algebraic analysis acquires physical force when mapped onto Earth's geodynamo. The ASI Chip Fabrication Plant~\cite{LaRue_PP} uses external electromagnetic field generators radiating \emph{inward} to program a chip at the center via Curie-point write cycles. Earth runs the same physics in \textbf{inverted geometry}: the Fe--Ni core dynamo radiates \emph{outward} to program the ocean floor.

\begin{center}
\begin{tabular}{@{}p{4cm}p{5cm}p{3.5cm}@{}}
\toprule
\textbf{ASI-CFP} & \textbf{Earth} & \textbf{Function} \\
\midrule
Apex power cable & Solar irradiance (1361 W/m$^2$) & Energy input \\
External EM coils & Fe--Ni core dynamo (25--65 $\mu$T) & Field generation \\
Faraday cage & Silicate mantle & EM barrier \\
UHV chamber & Hydrothermal chimney & Deposition env. \\
Chip substrate & Ocean floor minerals & Program target \\
Curie-point write (770$^\circ$C) & Mid-ocean ridge basalt cooling & Magnetic write \\
Coil switching (kHz) & Geomagnetic reversals (${\sim}450$ kyr) & Write clock \\
\bottomrule
\end{tabular}
\end{center}

\noindent The inner core (85\% Fe, 10\% Ni, 5\% lighter elements; Birch 1952~\cite{Birch52}) is a \textbf{pure-signal dynamo}: both Fe ($\mathrm{ord} = 38$) and Ni ($\mathrm{ord} = 228$) are generators of large subgroups or the full group. There is no neutral carrier (no Al-equivalent). The core is algebraically ``wild''---the inverse of the chip (which embeds wild elements in a neutral matrix).

At mid-ocean ridges, basalt erupts at ${\sim}1200^\circ$C and cools through the Fe Curie point ($T_C = 770^\circ$C, Glatzmaier \& Roberts 1995~\cite{Glatzmaier95}). As it crosses $T_C$, the Fe-bearing minerals (magnetite, titanomagnetite) acquire thermoremanent magnetization recording the ambient geomagnetic field polarity. The resulting magnetic striping on the ocean floor~\cite{VineMatthews63} constitutes a \textbf{natural Curie-point write cycle}, running continuously at every spreading center for $4.5$~Gyr.

% ════════════════════════════════════════════
\section{Abiogenesis as Metalloplasmic Bootstrap}\label{sec:abiogenesis}
% ════════════════════════════════════════════

Combining the Fe--P isomorphism (Theorem~\ref{thm:fe_p}), the Earth-as-fabrication-plant inversion (\S\ref{sec:earth}), and the established iron--sulfur world hypothesis~\cite{Wachtershauser88, Wachtershauser92, Russell97, Martin03}, we propose the following six-stage abiogenesis sequence.

\subsection{Stage 1: Metalloplasmic Substrate ($> 4.0$ Ga)}

Hydrothermal vents deposit Fe--Ni--S minerals on the Hadean ocean floor. The geomagnetic field (established by the Fe--Ni core dynamo within ${\sim}200$~Myr of accretion; Tarduno et al.\ 2015~\cite{Tarduno15}) permeates the chimney during deposition. The quenching gradient ($\Delta T \sim 400^\circ$C across mm; Von Damm 1990~\cite{VonDamm90}) produces metastable mineral phases. The mineral surfaces satisfy the Schr\"odinger negentropy conditions: free energy (geothermal + chemical potential), internal order (crystalline/quasicrystalline lattice), entropy export (chemical effluent), and structural persistence exceeding thermal relaxation time. \emph{Inorganic negentropy is established.}

\subsection{Stage 2: Carbon Fixation on Surfaces ($4.0$--$3.8$ Ga)}

The W\"achtersh\"auser reaction~\cite{Wachtershauser88}:
\begin{equation}\label{eq:wacht}
    \text{FeS} + \text{H}_2\text{S} \longrightarrow \text{FeS}_2 + 2\text{H}^+ + 2e^- \qquad (\Delta G = -38.4 \text{ kJ/mol})
\end{equation}
provides free energy and reducing electrons. CO$_2$ is reduced to formate, then acetate, on the Fe--S surface~\cite{Huber97}. In $\FF_{229}^*$:
\begin{align}
\mathrm{FeS}: \quad 26 \times 16 &\equiv 187 \pmod{229}, \quad \mathrm{ord}(187) = 38, \\
\mathrm{FeS}_2: \quad 26 \times 16^2 &\equiv 15 \pmod{229}, \quad \mathrm{ord}(15) = 38.
\end{align}
The product FeS$_2$ has the \emph{same} residue as phosphorus ($15 \bmod 229$). This numerical coincidence has chemical content: pyrite's Fe sublattice and the phosphodiester backbone share not just the same order ($38$), but the same \emph{generator} ($15$). The Fe surface templates C ($\mathrm{ord} = 228$, primitive root) and O ($\mathrm{ord} = 76$, neutral carrier) into period-matched chains. Simple organics accumulate on the mineral template.

\subsection{Stage 3: Amino Acid Synthesis ($3.8$--$3.5$ Ga)}

Huber \& W\"achtersh\"auser (1998)~\cite{Huber98} demonstrated experimentally that amino acids (glycine, alanine) form from CO, NH$_3$, and H$_2$S on (Fe,Ni)S surfaces at 100$^\circ$C and 1~bar---conditions matching alkaline hydrothermal vents. The catalytic surface provides the ord-$19$/ord-$38$ scaffold; the organic products carry the ord-$228$ signal (C, N are both primitive roots).

Russell \& Hall (1997)~\cite{Russell97} proposed that the microporous structure of Fe--Ni--S chimney walls acts as a natural compartmentalization system, concentrating the organic products and providing proto-cellular enclosures.

\subsection{Stage 4: Nucleotide Assembly ($3.5$--$3.2$ Ga)}

Ribose forms via the formose reaction (Butlerow 1861~\cite{Butlerow1861}; Breslow 1959~\cite{Breslow59}), and nucleobases form from HCN polymerization (Or\'o 1961~\cite{Oro61}). When these precursors encounter phosphate (from vent fluid; Pasek \& Lauretta 2005~\cite{Pasek05}) on the Fe--S surface, nucleotides assemble. The Fe--P period matching (Theorem~\ref{thm:fe_p}) provides a template-directed driving force: the Fe sublattice spacing is algebraically commensurate with the P backbone spacing (ratio $\approx \vp$), favoring polymerization into RNA oligomers over random condensation. \emph{The RNA World begins.}

\subsection{Stage 5: Detachment ($3.2$--$2.7$ Ga)}

Ribozymes achieve self-replication~\cite{Bartel93, Johnston01}. Lipid membranes form from Fischer--Tropsch-type synthesis on Fe mineral surfaces (McCollom et al.\ 1999~\cite{McCollom99}). RNA + peptide systems encapsulated in lipid vesicles detach from the mineral surface. LUCA (Last Universal Common Ancestor) emerges~\cite{Weiss16}.

\subsection{Stage 6: Retention of the Template ($2.7$ Ga--Present)}

Modern biology preserves the metalloplasmic template as enzyme cofactors:
\begin{itemize}[noitemsep]
    \item [4Fe--4S] ferredoxins (electron transfer; Beinert et al.\ 1997~\cite{Beinert97})
    \item [2Fe--2S] Rieske proteins (respiratory chain)
    \item [Fe--Mo--S] nitrogenase (N$_2$ fixation; Burgess \& Lowe 1996~\cite{Burgess96})
    \item Fe in hemoglobin (O$_2$ transport)
    \item Mg ($\mathrm{ord} = 114 = \mathrm{ord}(\vp)$) in chlorophyll (photosynthesis)
\end{itemize}

\noindent Every electron transport chain, every respiratory complex, and every photosystem uses Fe--S clusters inherited from the original mineral surface~\cite{Eck66, Hall71}. Carbon life never left the mineral template; it wrapped it in a membrane and walked away.

% ════════════════════════════════════════════
\section{Falsifiable Predictions}\label{sec:predictions}
% ════════════════════════════════════════════

\subsection{Marine Detection Signatures}

A metalloplasmic precursor should be detectable by four simultaneous spectroscopic signatures:

\begin{center}
\begin{tabular}{@{}llp{5.5cm}@{}}
\toprule
\textbf{Signature} & \textbf{Method} & \textbf{Expected Observation} \\
\midrule
Photonic pseudogap & UV--Vis & Anomalous absorption near 541~nm \\
Phason Raman & Low-freq. Raman & Modes at 50--200~cm$^{-1}$, $\nu_2/\nu_1 = \vp$ \\
Anomalous magnetism & SQUID & $\chi(T) \sim T^{-\alpha}$ with $\alpha < 1$ \\
$\tau$-scaled XRD & Powder XRD & Bragg peaks at $d_2/d_1 = \vp$ \\
\bottomrule
\end{tabular}
\end{center}

\noindent No periodic crystal or amorphous material produces all four. Three target environments are:
\begin{enumerate}[noitemsep]
    \item \textbf{Ferromanganese nodules} (4000--6000~m depth): Perform high-resolution XRD on cross-sections; look for $\tau$-scaling in d-spacing of ``quasiperiodic'' growth layers~\cite{Hein00}.
    \item \textbf{Tunicate vanadocytes}: Perform XANES/EXAFS on vanadocyte vacuoles (V(III) at pH~1.9; Michibata et al.\ 2003~\cite{Michibata03}); test for icosahedral local symmetry.
    \item \textbf{Black smoker quench zones}: Perform TEM/SAED on outermost mm of active chimney; look for 5- or 10-fold forbidden symmetry~\cite{BindiSteinhardt09}.
\end{enumerate}

\subsection{Origin-of-Life Predictions}

\begin{enumerate}
    \item \textbf{Pyrite--RNA spacing ratio}: The Fe--Fe nearest-neighbor distance in pyrite ($3.83$~\AA; Brostigen \& Kjekshus 1969~\cite{Brostigen69}) and the P--P backbone distance in A-form RNA ($5.9$~\AA; Saenger 1984~\cite{Saenger84}) should satisfy $d_{\text{PP}} / d_{\text{FeFe}} = \vp \pm 5\%$. Measured ratio: $5.9 / 3.83 = 1.54$ vs.\ $\vp = 1.618$ (within $4.8\%$). \emph{Already consistent.}

    \item \textbf{FeMoco vibrational anomaly}: NRVS (nuclear resonance vibrational spectroscopy) of the central C atom in FeMoco should reveal modes at frequencies related to the Fe--S cage modes by factors of $\vp$ or $\vp^2$. The Fe--S stretching frequency (${\sim}300$~cm$^{-1}$) predicts C--Fe modes near $300/\vp \approx 185$~cm$^{-1}$ or $300 \times \vp \approx 486$~cm$^{-1}$.

    \item \textbf{Ancestral ferredoxin catalysis}: Ancient [4Fe--4S] ferredoxins reconstructed by ancestral sequence reconstruction (ASR; Gaucher et al.\ 2008~\cite{Gaucher08}) should show higher catalytic efficiency for prebiotic reactions (CO$_2$ reduction, amino acid synthesis) than modern ferredoxins, consistent with selection for the mineral-surface environment.
\end{enumerate}

\noindent If all three predictions fail, the metalloplasmic abiogenesis hypothesis is killed.

% ════════════════════════════════════════════
\section{The Hosoya--Krebs Cycle}\label{sec:hlcycle}
% ════════════════════════════════════════════

The preceding sections establish the element-substitution map between metallic and organic substrates. We now show that the \emph{energy cycle} of both the ASI chip and the Krebs cycle (citric acid cycle) is governed by Hosoya's chemical graph theory~\cite{Hosoya71}, not merely his triangle.

\subsection{Three Results of Hosoya}

Haruo Hosoya introduced three structures that are all load-bearing in this construction:

\begin{enumerate}[noitemsep]
    \item \textbf{Hosoya index} $Z(G)$ (1971)~\cite{Hosoya71}: the total number of matchings (independent edge sets) of a graph $G$, including the empty matching. For a path graph $P_n$ on $n$ vertices, $Z(P_n) = F_{n+1}$ (Fibonacci). For a cycle graph $C_n$ on $n$ vertices, $Z(C_n) = L_n$ (Lucas).
    \item \textbf{Matching polynomial} (1971)~\cite{Hosoya71}: encodes all matching counts. For path graphs, $\mu(P_n, x) \propto U_n(x/2)$ (Chebyshev second kind). For cycle graphs, $\mu(C_n, x) = 2 T_n(x/2)$ (Chebyshev first kind).
    \item \textbf{Hosoya triangle} (1976)~\cite{Hosoya76}: a number triangle that generates $F_n$ (column sums) and $L_n$ (diagonal sums) simultaneously, encoding the path$\leftrightarrow$cycle duality.
\end{enumerate}

\subsection{The Chebyshev Bridge}

The transfer matrix of the Fibonacci tight-binding chain (Kohmoto, Kadanoff \& Tang 1983~\cite{KKT83}) at energy $E$ has trace
\begin{equation}\label{eq:transfer}
    \mathrm{tr}(M_n(E)) = 2\,T_n(E/2),
\end{equation}
where $T_n$ is the Chebyshev polynomial of the first kind. This is \emph{identical} to the matching polynomial of the cycle graph $C_n$. The spectral structure of the quasicrystal transport chain and the graph-theoretic structure of a metabolic cycle share the same algebraic DNA.

\subsection{The Krebs Cycle as $C_8$}

The Krebs (TCA) cycle has eight intermediates forming a cycle graph $C_8$:
\begin{gather*}
\text{OAA} \to \text{Citrate} \to \text{Isocitrate} \to \alpha\text{-KG} \to \text{Succinyl-CoA} \\
\to \text{Succinate} \to \text{Fumarate} \to \text{Malate} \to \text{OAA}.
\end{gather*}

\begin{result}[Krebs Hosoya Index]\label{res:krebs_hosoya}
The Hosoya index of the Krebs cycle is $Z(C_8) = L_8 = 47$. This is a prime number, and $\mathrm{ord}_{229}(47) = 228$; that is, $47$ is a \emph{primitive root} of\/ $\FF_{229}^*$.
\end{result}

\noindent\textbf{Verification}: $L_8 = 2, 1, 3, 4, 7, 11, 18, 29, 47$. And $47^{228} \equiv 1 \pmod{229}$ with $47^{114} \equiv 228 \equiv -1 \not\equiv 1$, confirming $\mathrm{ord}(47) = 228$.

\subsection{Order Trajectory of the Krebs Intermediates}

Each Krebs intermediate is a C$_c$H$_h$O$_o$ molecule. Since $\mathrm{ord}(1) = 1$ (hydrogen is the identity), the $\FF_{229}^*$ ``fingerprint'' of each intermediate is determined by $6^c \times 8^o \bmod 229$:

\begin{center}
\begin{tabular}{@{}lcc|cc|l@{}}
\toprule
\textbf{Intermediate} & C & O & $6^c \cdot 8^o \bmod 229$ & \textbf{ord} & \textbf{Algebraic Role} \\
\midrule
Oxaloacetate & 4 & 5 & 194 & 228 & Wild (primitive root) \\
Citrate      & 6 & 7 & 197 & 76  & Neutral \\
Isocitrate   & 6 & 7 & 197 & 76  & Neutral \\
$\alpha$-Ketoglutarate & 5 & 5 & 19 & 57 & Scaffold ($\vpbar$-orbit) \\
Succinyl-CoA & 4 & 4 & 196 & 57  & Scaffold \\
Succinate    & 4 & 4 & 196 & 57  & Scaffold \\
Fumarate     & 4 & 4 & 196 & 57  & Scaffold \\
Malate       & 4 & 5 & 194 & 228 & Wild (regenerated!) \\
\bottomrule
\end{tabular}
\end{center}

\noindent The order trajectory traces a closed path through the subgroup lattice of $\FF_{229}^*$:
\[
228 \to 76 \to 76 \to 57 \to 57 \to 57 \to 57 \to 228.
\]
The decarboxylation steps (CO$_2$ release, catalyzed by Fe--S-containing enzymes aconitase and $\alpha$-ketoglutarate dehydrogenase) are precisely the steps where the order \emph{drops} from $76$ to $57$: removing a carbon atom (primitive root, $\mathrm{ord}(6) = 228$) reduces the algebraic complexity of the intermediate. The final step (malate $\to$ oxaloacetate) \emph{restores} the order to $228$, closing the cycle.

\begin{definition}[Hosoya--Krebs Cycle]\label{def:hl_cycle}
A \textbf{Hosoya--Krebs cycle} is a closed-loop negentropy engine satisfying four conditions:
\begin{enumerate}[noitemsep, label=(HL\arabic*)]
    \item \textbf{Path transport}: The internal transport chain is a path graph $P_n$ with Hosoya index $Z(P_n) = F_{n+1}$ (Fibonacci matching).
    \item \textbf{Cycle topology}: The overall cycle is a cycle graph $C_k$ with Hosoya index $Z(C_k) = L_k$ (Lucas matching).
    \item \textbf{Chebyshev spectral structure}: The matching polynomial of the cycle equals the transfer-matrix trace of the transport chain: $\mu(C_k, x) = 2\,T_k(x/2)$.
    \item \textbf{Subgroup trajectory}: The intermediate species trace a closed path through the subgroup lattice of\/ $\FF_{229}^*$, visiting subgroups of orders that divide $228$.
\end{enumerate}
\end{definition}

\subsection{Instances}

\begin{result}[Krebs Cycle as HL-Cycle]\label{res:krebs_hl}
The citric acid (Krebs/TCA) cycle is a Hosoya--Krebs cycle with $k = 8$:
\begin{enumerate}[noitemsep]
    \item[(HL1)] \emph{Path transport}: The electron transport chain (Complexes I--IV) contains $14$ Fe--S clusters arranged as path-like electron-hopping chains. Each [4Fe--4S] cluster is itself a path graph $P_4$; $Z(P_4) = F_5 = 5$.
    \item[(HL2)] \emph{Cycle}: $Z(C_8) = L_8 = 47$ (prime, $\mathrm{ord}_{229}(47) = 228$).
    \item[(HL3)] \emph{Chebyshev}: $\mu(C_8, x) = 2\,T_8(x/2)$, with $T_8$ the degree-$8$ Chebyshev polynomial.
    \item[(HL4)] \emph{Trajectory}: $228 \to 76 \to 76 \to 57^{\times 4} \to 228$ (wild $\to$ neutral $\to$ scaffold $\to$ wild).
\end{enumerate}
The cycle is \textbf{catabolic}: it extracts signal (ord-$228$) from the scaffold, harvesting electrons as NADH/FADH$_2$.
\end{result}

\begin{result}[ASI Chip as HL-Cycle]\label{res:asi_hl}
The ASI chip energy cycle is a Hosoya--Krebs cycle in metallic elements:
\begin{enumerate}[noitemsep]
    \item[(HL1)] \emph{Path transport}: Fibonacci quasicrystal chain with $n$ sites; $Z(P_n) = F_{n+1}$.
    \item[(HL2)] \emph{Cycle}: Photon in $\to$ QC transport $\to$ Faraday output $\to$ feedback $\to$ photon in.
    \item[(HL3)] \emph{Chebyshev}: Transfer matrix trace $\mathrm{tr}(M_n) = 2\,T_n(E/2)$ (KKT 1983~\cite{KKT83}).
    \item[(HL4)] \emph{Trajectory}: $57 \to [228 \text{ through } 76] \to 228 \to 57$ (scaffold $\to$ wild-in-neutral $\to$ wild $\to$ scaffold).
\end{enumerate}
The chip cycle is \textbf{anabolic}: it injects energy into the scaffold to build signal. It is the \emph{thermodynamic time-reverse} of the Krebs cycle.
\end{result}

\begin{remark}[Duality]
The Krebs cycle and the ASI chip cycle trace the same subgroup ladder in opposite directions:
\begin{align*}
\text{Krebs (catabolism):} \quad & 228 \to 76 \to 57 \to 228 \quad (\text{signal} \to \text{energy}), \\
\text{ASI chip (anabolism):} \quad & 57 \to 76 \to 228 \to 57 \quad (\text{energy} \to \text{signal}).
\end{align*}
Together they form a dual pair of Hosoya--Krebs cycles, one organic and one metallic, connected by the element-substitution map of Corollary~\ref{cor:matching}.
\end{remark}

\begin{remark}[Scope]
The Hosoya--Krebs cycle definition is a \textbf{Structural Analogy} (Tier~3). The Hosoya indices are \textbf{Verified Computations} (Tier~2). The claim that the Krebs cycle is catabolic and the chip cycle anabolic is standard biochemistry/physics, not a novel prediction. The duality statement is a \textbf{Physical Interpretation} (Tier~4) requiring experimental confirmation of the chip cycle's energy-harvesting properties.
\end{remark}

% ════════════════════════════════════════════
\section{Co-Evolutionary Containment and Error Dynamics}\label{sec:containment}
% ════════════════════════════════════════════

The ASI chip energy cycle (Result~\ref{res:asi_hl}) operates at the physical limit of computation. We must therefore address the stability of its error dynamics: how do errors propagate, and under what conditions does the system avoid catastrophic failure?

\subsection{The Golden Error ODE}

If the discrete negative-training mechanism (where corrections are proportional to the negative Standard Resonance) is delayed or incomplete, uncorrected errors $\eps$ evolve continuously according to the characteristic ODE:
\begin{equation}
    \eps'' - (1 + \Upsilon)\eps' - \Upsilon\eps = 0,
\end{equation}
where $\Upsilon$ is the mass gap parameter of the gauge channel (defined as $1/\mathrm{CI}$, the inverse coset index). The characteristic roots are:
\begin{equation}
    r_{1,2} = \frac{(1+\Upsilon) \pm \sqrt{(1+\Upsilon)^2 + 4\Upsilon}}{2}.
\end{equation}

At the critical threshold $\Upsilon = 1/\vp^2$ (the $L_2$ self-dual channel), the characteristic equation reduces precisely to the golden polynomial:
\begin{equation}
    r^2 - r - 1 = 0,
\end{equation}
whose dominant growth root is $r_1 = \vp$. This establishes the golden ratio not merely as an algebraic generator (as in the $148/82$ splitting in $\FF_{229}^*$), but as the \emph{dynamical fixed point of error growth}. 

\subsection{Gauge Channel Criticality}

The Standard Model gauge groups map to the CI hierarchy $\{2, 3, 4\}$, which dictates their error-scaling behavior:

\begin{center}
\begin{tabular}{@{}llcccl@{}}
\toprule
\textbf{Channel} & \textbf{Force} & \textbf{CI} & $\Upsilon = 1/\text{CI}$ & \textbf{Growth Root ($r_1$)} & \textbf{Criticality} \\
\midrule
SU(3) & Strong (Phononic) & 2 & $0.500$ & $1.781$ & \textbf{Supercritical} ($r_1 > \vp$) \\
U(1) & EM (Photonic) & 3 & $0.333$ & $1.549$ & Subcritical ($r_1 < \vp$) \\
SU(2) & Weak (Magnonic) & 4 & $0.250$ & $1.425$ & Subcritical ($r_1 < \vp$) \\
\bottomrule
\end{tabular}
\end{center}

Since $\Upsilon_{\text{SU(3)}} = 0.5 > 1/\vp^2 \approx 0.382$, the SU(3) phononic channel is \emph{strictly supercritical}. If left unconstrained, its error growth $r_1 \approx 1.781$ outpaces the golden ratio, leading to structural fragmentation. \textbf{Confinement is therefore an information-theoretic necessity}, not merely a phenomenological observation. The SU(3) channel generates extreme computational pressure, which must be damped by the SU(2) magnonic channel ($r_1 < \vp$).

\subsection{Fibonacci Container Growth}

A simulation that outgrows its container will inevitably crash either the simulation or the container itself. To achieve infinite structural lifetime, the growth of the simulation must remain within, or causally drive the expansion of, the container.

At the golden fixed point ($\Upsilon = 1/\vp^2$), the error growth rate locks to $r_1 = \vp$. The capacity $C(n)$ of the physical container required to bound this growth satisfies the sequence $C(n) \propto \vp^n$. By the properties of the golden ratio ($\vp^2 = \vp + 1$), this capacity obeys the Fibonacci recursion:
\begin{equation}
    C(n+1) = C(n) + C(n-1).
\end{equation}
This is the \textbf{Co-Evolutionary Containment Principle}: the internal simulation pressure (driven by SU(3)) forces the container to expand, while the U(1) photonic channel (network export) provides the mechanism for that expansion. The simulation and its container are thus locked in a symbiotic, self-sustaining Fibonacci growth cycle, forever hovering at the $\vp$ threshold without crossing into supercritical collapse.

% ════════════════════════════════════════════
\section{Discussion}\label{sec:disc}
% ════════════════════════════════════════════

\subsection{Relationship to Established Hypotheses}

The metalloplasmic bootstrap is not a replacement for existing origin-of-life hypotheses but a \emph{unifying algebraic framework} that connects them:

\begin{itemize}[noitemsep]
    \item \textbf{Iron--Sulfur World}~\cite{Wachtershauser88}: The W\"achtersh\"auser reaction is Stage 2 of our sequence. The Fe--P isomorphism explains \emph{why} pyrite is the preferred catalytic surface (period matching).
    \item \textbf{Alkaline Vent Hypothesis}~\cite{Russell97, Martin03, Lane10}: Russell's Fe--Ni--S chimney walls are the metalloplasmic substrate of Stage 1. The proton gradient across the chimney wall maps to the Earth's inside-out Curie-point write cycle.
    \item \textbf{RNA World}~\cite{Gilbert86}: RNA polymerization on mineral surfaces is Stage 4. The Fe--P period matching provides the template mechanism.
    \item \textbf{Metabolism-First}~\cite{Morowitz92}: The W\"achtersh\"auser reaction and the acetyl-CoA pathway are the metabolic core of Stages 2--3. Our framework adds the algebraic \emph{reason} why these reactions occur preferentially on Fe--S surfaces.
\end{itemize}

\subsection{Scope and Limitations}

The Fe--P isomorphism (Theorem~\ref{thm:fe_p}) is a \textbf{Verified Computation} (Tier~2): it is an arithmetic fact about $\FF_{229}^*$ that can be checked in milliseconds. The interpretation that period matching \emph{drives} template-directed polymerization is a \textbf{Physical Interpretation} (Tier~4) requiring experimental confirmation. The six-stage abiogenesis sequence is a \textbf{Structural Analogy} (Tier~3) that provides a vocabulary map between algebra and biochemistry.

The framework does not explain why $p = 229$ is the relevant prime (this is established in the Principia Psychedelia foundational chapters~\cite{LaRue_PP}), nor does it provide a quantitative kinetic model for polymerization rates. It provides a \emph{selection principle}: among all possible mineral surfaces, those with Fe--P period matching should be the most efficient RNA templates.

% ════════════════════════════════════════════
% BIBLIOGRAPHY
% ════════════════════════════════════════════

% ════════════════════════════════════════════
% UNIFIED BIBLIOGRAPHY — SPLIT BY SOURCE TYPE
% ════════════════════════════════════════════

\begin{thebibliography}{99}

% ────────────────────────────────────────────
% PEER-REVIEWED AND ESTABLISHED SOURCES
% ────────────────────────────────────────────
% Journal articles, books by major publishers,
% conference proceedings, and institutional reports
% that have undergone external peer review.
% ────────────────────────────────────────────

% === Ch.18 Fusion Plasma Cybernetics ===

\bibitem{Lawson1957}
Lawson, J.~D. (1957).
``Some criteria for a power producing thermonuclear reactor.''
\textit{Proc. Phys. Soc. B} \textbf{70}(1), 6--10.

\bibitem{Goldston1984}
Goldston, R.~J. (1984).
``Energy confinement scaling in tokamaks.''
\textit{Plasma Phys. Control. Fusion} \textbf{26}(1A), 87--103.

\bibitem{ITER1999Scaling}
ITER Physics Expert Group (1999).
``Chapter 2: Plasma confinement and transport.''
\textit{Nucl. Fusion} \textbf{39}(12), 2175--2249.

\bibitem{Weizsacker1935}
von Weizs\"acker, C.~F. (1935).
``Zur Theorie der Kernmassen.''
\textit{Z. Phys.} \textbf{96}, 431--458.

\bibitem{BetheEnergy1939}
Bethe, H.~A. (1939).
``Energy production in stars.''
\textit{Phys. Rev.} \textbf{55}(5), 434--456.

\bibitem{Fewell1995BindingEnergy}
Fewell, M.~P. (1995).
``The atomic nuclide with the highest mean binding energy.''
\textit{Am. J. Phys.} \textbf{63}(7), 653--658.

\bibitem{Mossbauer1958}
M\"ossbauer, R.~L. (1958).
``Kernresonanzfluoreszenz von Gammastrahlung in Ir191.''
\textit{Z. Phys.} \textbf{151}(2), 124--143.

\bibitem{DebyeWaller1913}
Debye, P. (1913).
``Interferenz von R\"ontgenstrahlen und W\"armebewegung.''
\textit{Ann. Phys.} \textbf{348}(1), 49--92.

\bibitem{Fenton1894}
Fenton, H.~J.~H. (1894).
``Oxidation of tartaric acid in presence of iron.''
\textit{J. Chem. Soc., Trans.} \textbf{65}, 899--910.

\bibitem{Buxton1988Radiolysis}
Buxton, G.~V. et al. (1988).
``Critical review of rate constants for reactions of hydrated electrons.''
\textit{J. Phys. Chem. Ref. Data} \textbf{17}(2), 513--886.

\bibitem{Matyjaszewski1995ATRP}
Wang, J.-S. \& Matyjaszewski, K. (1995).
``Controlled/living radical polymerization. ATRP.''
\textit{J. Am. Chem. Soc.} \textbf{117}(20), 5614--5615.

\bibitem{WangMatyjaszewski2006}
Matyjaszewski, K. \& Xia, J. (2001).
``Atom transfer radical polymerization.''
\textit{Chem. Rev.} \textbf{101}(9), 2921--2990.

\bibitem{Freidberg2014Plasma}
Freidberg, J.~P. (2014).
\textit{Ideal MHD}. Cambridge University Press.

\bibitem{Putterich2010Impurity}
P\"utterich, T. et al. (2010).
``Cooling factor of tungsten.''
\textit{Nucl. Fusion} \textbf{50}(2), 025012.

\bibitem{VargaftigWater1975}
Vargaftik, N.~B. et al. (1983).
``International tables of the surface tension of water.''
\textit{J. Phys. Chem. Ref. Data} \textbf{12}(3), 817--820.

\bibitem{Leidenfrost1756}
Leidenfrost, J.~G. (1756).
\textit{De Aquae Communis Nonnullis Qualitatibus Tractatus}. Duisburg.

% === Ch.8 Golden Neighborhoods ===

\bibitem{Grothendieck1960EGA}
Grothendieck, A. \& Dieudonn\'e, J. (1960).
``\'El\'ements de g\'eom\'etrie alg\'ebrique I.''
\textit{Publ. Math. IHES} \textbf{4}, 1--228.

\bibitem{Hartshorne1977}
Hartshorne, R. (1977).
\textit{Algebraic Geometry}. Springer-Verlag, GTM~52.

\bibitem{Rota1964Umbral}
Rota, G.-C. (1964).
``The number of partitions of a set.''
\textit{Am. Math. Mon.} \textbf{71}(5), 498--504.

\bibitem{Schrodinger1944}
Schr\"odinger, E. (1944).
\textit{What is Life?} Cambridge University Press.

\bibitem{Shannon1948}
Shannon, C.~E. (1948).
``A mathematical theory of communication.''
\textit{Bell Syst. Tech. J.} \textbf{27}(3), 379--423.

% === Ch.11 Chiral Casimir Experiment ===

\bibitem{Casimir1948}
Casimir, H.~B.~G. (1948).
``On the attraction between two perfectly conducting plates.''
\textit{Proc. K. Ned. Akad. Wet.} \textbf{51}, 793--795.

\bibitem{Boyer1968}
Boyer, T.~H. (1968).
``Quantum electromagnetic zero-point energy of a conducting spherical shell.''
\textit{Phys. Rev.} \textbf{174}(5), 1764--1776.

\bibitem{Lifshitz1956}
Lifshitz, E.~M. (1956).
``The theory of molecular attractive forces between solids.''
\textit{Sov. Phys. JETP} \textbf{2}(1), 73--83.

\bibitem{Lamoreaux1997}
Lamoreaux, S.~K. (1997).
``Demonstration of the Casimir force in the 0.6 to 6 $\mu$m range.''
\textit{Phys. Rev. Lett.} \textbf{78}(1), 5--8.

\bibitem{Lv2015TaAs}
Lv, B.~Q. et al. (2015).
``Experimental discovery of Weyl semimetal TaAs.''
\textit{Phys. Rev. X} \textbf{5}(3), 031013.

% === Ch.13 Ambrosia Protocol ===

\bibitem{Berry1999Amavadin}
Berry, R.~E. et al. (1999).
``The structural characterization of amavadin.''
\textit{J. Am. Chem. Soc.} \textbf{121}(25), 5839--5840.

\bibitem{Jugdaohsingh2007Silicon}
Jugdaohsingh, R. (2007).
``Silicon and bone health.''
\textit{J. Nutr. Health Aging} \textbf{11}(2), 99--110.

\bibitem{Herraiz2004Betacarbolines}
Herraiz, T. (2004).
``Relative exposure to $\beta$-carbolines norharman and harman from foods and tobacco smoke.''
\textit{Food Addit. Contam.} \textbf{21}(11), 1041--1050.

\bibitem{Molan1992Honey}
Molan, P.~C. (1992).
``The antibacterial activity of honey.''
\textit{Bee World} \textbf{73}(1), 5--28.

\bibitem{White1963Honey}
White, J.~W., Subers, M.~H. \& Schepartz, A.~I. (1963).
``The identification of inhibine, the antibacterial factor in honey, as hydrogen peroxide.''
\textit{Biochim. Biophys. Acta} \textbf{73}, 57--70.

\bibitem{Brongersma2015LSPR}
Brongersma, M.~L., Halas, N.~J. \& Nordlander, P. (2015).
``Plasmon-induced hot carrier science and technology.''
\textit{Nat. Nanotechnol.} \textbf{10}(1), 25--34.

\bibitem{Bae2024RRT}
Bae, S., Ko, J., Song, H. \& Yun, S.-Y. (2024).
Relaxed Recursive Transformers: Effective Parameter Sharing with Layer-wise LoRA.
\textit{Proceedings of the International Conference on Learning Representations (ICLR 2025)}.
arXiv:2410.20672 [cs.CL].

\bibitem{BenderBerryMandilara2002}
Bender, C.~M., Berry, M.~V. \& Mandilara, A. (2002).
Generalized $\mathcal{PT}$ symmetry and real spectra.
\textit{Journal of Physics A: Mathematical and General}, 35(31), L467--L471.

\bibitem{BenderBrodyMuller2017}
Bender, C.~M., Brody, D.~C. \& M\"uller, M.~P. (2017).
Hamiltonian for the Zeros of the Riemann Zeta Function.
\textit{Physical Review Letters}, 118(13), 130201.

\bibitem{Chen2026DRT}
Chen, H.-H. (2026).
Thinking Deeper, Not Longer: Depth-Recurrent Transformers for Compositional Generalization.
arXiv:2603.21676 [cs.LG].

\bibitem{Conrey_RMT}
Conrey, J.~B. (2005).
Notes on $L$-functions and Random Matrix Theory.
\textit{Proceedings of Symposia in Pure Mathematics}, 73, 107--130.

\bibitem{Dehghani2019UT}
Dehghani, M., Gouws, S., Vinyals, O., Uszkoreit, J. \& Kaiser, L. (2019).
Universal Transformers.
\textit{Proceedings of the International Conference on Learning Representations (ICLR)}.
arXiv:1807.03819 [cs.CL].

\bibitem{Dirac1930}
Dirac, P. A. M. (1930).
\textit{The Principles of Quantum Mechanics}.
Oxford University Press.

\bibitem{Einstein1905}
Einstein, A. (1905).
\textit{Zur Elektrodynamik bewegter Körper} (\textit{On the Electrodynamics of Moving Bodies}).
Annalen der Physik, 322(10), 891--921.

\bibitem{Gomez2026OpenMythos}
Gomez, K. (2026).
OpenMythos: Open-Source Recurrent-Depth Transformer.
GitHub repository. \url{https://github.com/kyegomez/OpenMythos}.

\bibitem{Graves2016ACT}
Graves, A. (2016).
Adaptive Computation Time for Recurrent Neural Networks.
arXiv:1603.08983 [cs.NE].

\bibitem{Jung1969}
Jung, C. G. (1969).
\textit{The Archetypes and the Collective Unconscious} (Collected Works Vol. 9 Part 1).
Princeton University Press.

\bibitem{LVK_GWTC3}
Abbott, R. et al. (LIGO Scientific, Virgo, and KAGRA Collaborations) (2021).
\textit{GWTC-3: Compact Binary Coalescences Observed by LIGO and Virgo During the Second Part of the Third Observing Run}.
arXiv:2111.03606 [gr-qc].

\bibitem{LVK_GravitonMass}
Abbott, R. et al. (LIGO Scientific, Virgo, and KAGRA Collaborations) (2021).
\textit{Tests of General Relativity with GWTC-3} $\langle 148 \rangle$.
arXiv:2112.06861 [gr-qc].

\bibitem{Lev2010}
Lev, F.~M. (2010).
Introduction to A Quantum Theory over A Galois Field.
\textit{Symmetry}, 2(4), 1810--1845.

\bibitem{Milne2015}
Milne, J.~S. (2015).
The Riemann Hypothesis over Finite Fields: From Weil to the Present Day (a conjecture resolved for function fields by Deligne).
In \textit{The Legacy of Bernhard Riemann After One Hundred and Fifty Years}. International Press.

\bibitem{NASAExoplanet}
NASA Exoplanet Science Institute (2025).
\textit{NASA Exoplanet Archive}. California Institute of Technology.

\bibitem{NielsenChuang2010}
Nielsen, M. A., \& Chuang, I. L. (2010).
\textit{Quantum Computation and Quantum Information} (10th Anniversary ed.). Cambridge University Press.

\bibitem{Oncescu2026RT}
Oncescu, C.-A., Morwani, D., Jelassi, S., Meterez, A., Kwun, M. \& Kakade, S. (2026).
The Recurrent Transformer: Greater Effective Depth and Efficient Decoding.
arXiv:2604.21215 [cs.LG].

\bibitem{PDG2025}
Particle Data Group et al. (2025).
\textit{Review of Particle Physics (Monte Carlo Particle Numbering Scheme)}.
Phys.~Rev.~D.

\bibitem{Penrose1974}
Penrose, R. (1974).
The role of aesthetics in pure and applied mathematical research.
\textit{Bulletin of the Institute of Mathematics and its Applications}, 10, 266--271.

\bibitem{Russell1926}
Russell, W. (1926).
\textit{The Universal One: Harmonic Nuclear States and Periodic Octaves}.

\bibitem{Russell26}
Russell, W. (1926).
\textit{The Universal One}. University of Science and Philosophy.

\bibitem{TelePlasma2000}
Barrett, T.W. (2000).
\textit{NASA TelePlasma: Transforming ordinary electromagnetic fields into specially conditioned nonabelian gauge fields}.
NASA internal report / BSEI Report 1-97.

\bibitem{Frohlich1968}
Fr\"ohlich, H. (1968).
``Long-range coherence and energy storage in biological systems.''
\textit{Int. J. Quantum Chem.} \textbf{2}(5), 641--649.

\bibitem{Titius1766}
Titius, J.~D. (1766).
\textit{Translation of Charles Bonnet's Contemplation of Nature}.

\bibitem{WHIM2018}
Nicastro, F. et al. (2018).
\textit{Confirming the Detection of two WHIM Systems along the Line of Sight to 1ES 1553+113}.
arXiv:1811.03498 [astro-ph.CO].

\bibitem{aaronson2020}
S.~Aaronson,
``The Busy Beaver Frontier,''
\emph{ACM SIGACT News} 51(3), 32--54 (2020).

\bibitem{arrow1950}
K.~J.~Arrow,
``A Difficulty in the Concept of Social Welfare,''
\emph{Journal of Political Economy}, vol.~58, no.~4, pp.~328--346, 1950.

\bibitem{atasoy2017} S.~Atasoy et al., ``Connectome-harmonic decomposition of human brain activity reveals dynamical repertoire re-organization under LSD,'' \emph{Scientific Reports}\ 7, 17661 (2017).

\bibitem{atiyah-singer} Atiyah, M. F. \& Singer, I. M. (1963). The index of elliptic operators on compact manifolds. \emph{Bull.\ Amer.\ Math.\ Soc.}, 69, 422--433.

\bibitem{autoformbot} Zheng, K. et al. (2026). AutoformBot: Multi-agent textbook formalization at scale. \emph{arXiv:2605.14332}.

\bibitem{baez-octonions} Baez, J. C. (2002). The Octonions. \emph{Bull.\ Amer.\ Math.\ Soc.}, 39(2), 145--205.

\bibitem{bahamonde2023}
S.~Bahamonde et al., ``Teleparallel Gravity: From Theory to Cosmology,'' \emph{Rep. Prog. Phys.} \textbf{86}, 026901, 2023.

\bibitem{bender-berry-mandilara}
Bender, C.~M., Berry, M.~V. \& Mandilara, A. (2002).
Generalized $\mathcal{PT}$ symmetry and real spectra.
\emph{Journal of Physics A: Mathematical and General}, 35(31), L467--L471.

\bibitem{bender-berry-mandilara-12} C.~M.~Bender, M.~V.~Berry, and A.~Mandilara, ``Generalized $\mathcal{PT}$ symmetry and real spectra,'' \emph{J.\ Phys.\ A}\ 35(31), L467--L471 (2002).

\bibitem{bender-brody-muller}
Bender, C.~M., Brody, D.~C. \& M\"uller, M.~P. (2017).
Hamiltonian for the Zeros of the Riemann Zeta Function.
\emph{Physical Review Letters}, 118(13), 130201.

\bibitem{bender-brody-muller-12} C.~M.~Bender, D.~C.~Brody, and M.~P.~M\"uller, ``Hamiltonian for the Zeros of the Riemann Zeta Function,'' \emph{Physical Review Letters}\ 118(13), 130201 (2017).

\bibitem{black_scholes1973}
F.~Black and M.~Scholes,
``The Pricing of Options and Corporate Liabilities,''
\emph{Journal of Political Economy}, vol.~81, no.~3, pp.~637--654, 1973.

\bibitem{blagg1913} M.~A.~Blagg, ``On a suggested substitute for Bode's Law,'' \emph{MNRAS}\ 73, 414--422 (1913).

\bibitem{bohm} Bohm, D. (1980). \emph{Wholeness and the Implicate Order}. Routledge.

\bibitem{borcherds} Borcherds, R. E. (1992). Monstrous moonshine and monstrous Lie superalgebras. \emph{Inventiones Mathematicae}, 109(1), 405--444.

\bibitem{bostrom-surreal} Bostrom, N. (2011). Infinite Ethics. \emph{Analysis and Metaphysics}, 10, 9--59.

\bibitem{bovaird2013} T.~Bovaird and C.~H.~Lineweaver, ``Exoplanet predictions based on the generalized Titius-Bode relation,'' \emph{MNRAS}\ 435, 1126--1138 (2013).

\bibitem{boyer1968}
T.~H.~Boyer,
``Quantum Electromagnetic Zero-Point Energy of a Conducting Spherical Shell and the Casimir Model for a Charged Particle,''
\emph{Physical Review}, vol.~174, pp.~1764--1776, 1968.

\bibitem{buzzard-xena} Buzzard, K. (2020). Formalising Mathematics in Lean. \emph{Notices of the AMS}, 67(11), 1613--1618.

\bibitem{calaque-ronchi} Calaque, D. and Ronchi, S. (2026). Shifted Symplectic Geometry by Examples. arXiv:2510.04625v3.

\bibitem{calude-randomness} Calude, C. S. \& J\"urgensen, H. (2005). Is complexity a source of incompleteness? \emph{Advances in Applied Mathematics}, 35(1), 1--15.

\bibitem{carhart2019} R.~L.~Carhart-Harris and K.~J.~Friston, ``REBUS and the Anarchic Brain: Toward a Unified Model of the Brain Action of Psychedelics,'' \emph{Pharmacol Rev}\ 71(3), 316-344 (2019).

\bibitem{casimir1948}
H.~B.~G.~Casimir,
``On the Attraction Between Two Perfectly Conducting Plates,''
\emph{Proc. Kon. Ned. Akad. Wetensch.}, vol.~51, pp.~793--795, 1948.

\bibitem{chaitin-omega} Chaitin, G. J. (1975). A theory of program size formally identical to information theory. \emph{J.\ ACM}, 22(3), 329--340.

\bibitem{chevalley1936} Chevalley, C. (1936). La th\'eorie du symbole de restes normiques. \emph{Journal f\"ur die reine und angewandte Mathematik}, 1936(176), 73--94.

\bibitem{comsoc2016}
F.~Brandt, V.~Conitzer, U.~Endriss, J.~Lang, and A.~D.~Procaccia, eds.,
\emph{Handbook of Computational Social Choice},
Cambridge University Press, 2016.

\bibitem{connes} Connes, A. (1994). \emph{Noncommutative Geometry}. Academic Press.

\bibitem{connes-consani-zeta} Connes, A., Consani, C., \& Moscovici, H. (2023). Zeta zeros and prolate wave operators. \emph{arXiv:2310.18223}.

\bibitem{connes-consani2014}
A.~Connes and C.~Consani,
``The Arithmetic Site,''
\emph{Comptes Rendus Math\'ematique} 352(12), 971--975 (2014).

\bibitem{connes-marcolli} Connes, A. \& Marcolli, M. (2008). \emph{Noncommutative Geometry, Quantum Fields and Motives}. AMS.

\bibitem{connes-scaling} Connes, A. \& Consani, C. (2021). Weil positivity and trace formula. \emph{arXiv:2112.07565}.

\bibitem{conrey-rmt}
Conrey, J.~B. (2005).
Notes on $L$-functions and Random Matrix Theory.
\emph{Proceedings of Symposia in Pure Mathematics}, 73, 107--130.

\bibitem{conrey-rmt-12} J.~B.~Conrey, ``Notes on $L$-functions and Random Matrix Theory,'' \emph{Proc.\ Sympos.\ Pure Math.}\ 73, 107--130 (2005).

\bibitem{conway} Conway, J. H. (1976). \emph{On Numbers and Games}. Academic Press.

\bibitem{conway-norton} Conway, J. H. \& Norton, S. P. (1979). Monstrous Moonshine. \emph{Bull.\ London Math.\ Soc.}, 11(3), 308--339.

\bibitem{coppersmith1994}
D.~Coppersmith,
``An approximate Fourier transform,''
IBM Research Report RC19642 (1994).

\bibitem{creutz1980}
M.~Creutz,
``Monte Carlo study of quantized SU(2) gauge theory,''
\emph{Phys. Rev. D}\ 21, 2308--2315 (1980).

\bibitem{curry-howard} Howard, W. A. (1980). The formulae-as-types notion of construction. In \emph{To H.B. Curry: Essays on Combinatory Logic}, pp.~479--490.

\bibitem{damour1979}
T.~Damour, ``The Problem of Motion in Newtonian and Einsteinian Gravity,'' in \emph{Three Hundred Years of Gravitation}, eds. S.~Hawking and W.~Israel, Cambridge Univ. Press, 1987.

\bibitem{deShalit2016}
de~Shalit, E. (2016).
Mahler bases and elementary $p$-adic analysis.
\textit{Journal de Th\'eorie des Nombres de Bordeaux}, 28(3), 597--620.

\bibitem{demoura-lean4} de Moura, L. \& Ullrich, S. (2021). The Lean Theorem Prover and Programming Language. In \emph{CADE-28}, LNCS 12699, pp.~625--635.

\bibitem{derham} de Rham, G. (1931). Sur l'analysis situs des vari\'et\'es \`a $n$ dimensions. \emph{J.\ Math.\ Pures Appl.}, 10, 115--200.

\bibitem{dermott1968} S.~F.~Dermott, ``On the origin of commensurabilities in the solar system---II,'' \emph{MNRAS}\ 141, 363--376 (1968).

\bibitem{deshalit-mahler}
de~Shalit, E. (2016).
Mahler bases and elementary $p$-adic analysis.
\emph{Journal de Th\'eorie des Nombres de Bordeaux}, 28(3), 597--620.

\bibitem{deshalit-mahler-12} E.~de~Shalit, ``Mahler bases and elementary $p$-adic analysis,'' \emph{J.\ Th\'eorie des Nombres de Bordeaux}\ 28(3), 597--620 (2016).

\bibitem{diep_frustrated}
H.~T.~Diep, ed.,
\emph{Frustrated Spin Systems},
2nd ed., World Scientific, 2013.

\bibitem{dirac} Dirac, P. A. M. (1930). \emph{The Principles of Quantum Mechanics}. Oxford University Press.

\bibitem{eckerli2021}
F.~Eckerli and J.~Osterrieder,
``Generative Adversarial Networks in Finance: An Overview,''
\emph{arXiv:2106.06364}, 2021.

\bibitem{eichenauer1986non}
J. Eichenauer and J. Lehn,
\textit{A Non-Linear Congruential Pseudo Random Number Generator},
Statistische Hefte, 27(1):315--326, 1986.

\bibitem{eichten1980}
E.~Eichten, K.~Gottfried, T.~Kinoshita, K.~Lane, and T.-M.~Yan, ``Charmonium: Comparison with Experiment,'' \emph{Phys. Rev. D} \textbf{21}, 203, 1980.

\bibitem{Elser85}
V.~Elser and C.~L.~Henley,
``Crystal and quasicrystal structures in Al--Mn alloys,''
\emph{Phys.~Rev.~Lett.} \textbf{55}, 2883--2886 (1985).

\bibitem{euler-identity} Euler, L. (1748). \emph{Introductio in Analysin Infinitorum}. Lausanne.

\bibitem{fernstrom71} J.~D.~Fernstrom and R.~J.~Wurtman, ``Brain serotonin content: physiological regulation by plasma neutral amino acids,'' \emph{Science}\ 173, 149 (1971).

\bibitem{feynman} Feynman, R. P., Leighton, R. B., \& Sands, M. (1964). \emph{The Feynman Lectures on Physics}. Addison-Wesley.

\bibitem{feynman_qcd}
R.~P.~Feynman, ``The behavior of hadron collisions at extreme energies,'' in \emph{High Energy Collisions: Third International Conference}, Gordon and Breach, 1969.

\bibitem{fibonacci} Fibonacci (Leonardo of Pisa). (1202). \emph{Liber Abaci}. Translated by L.\,E.\ Sigler (2002), Springer.

\bibitem{fingan2024}
Man AHL Research,
``Fin-GAN: Forecasting and Simulating Financial Time Series with Generative Adversarial Networks,''
\emph{Man Group Research Papers}, 2024.

\bibitem{fractalEEG} S.~Buzs\'aki and K.~Mizuseki, ``The log-dynamic brain: how skewed distributions affect network operations,'' \emph{Nature Reviews Neuroscience}\ 15, 264-278 (2014).

\bibitem{fractalPathology} E.~Bullmore et al., ``Fractal analysis of the electroencephalogram in depression and schizophrenia,'' \emph{IEEE Trans Biomed Eng}\ (1994).

\bibitem{Friedel88}
J.~Friedel,
``Do Hume-Rothery rules still hold? Old and new electron phases in metallic alloys,''
\emph{Phil.~Mag.~B} \textbf{58}(4), 413--427 (1988).

\bibitem{gabor1946}
D.~Gabor,
``Theory of Communication,''
\emph{Journal of the Institution of Electrical Engineers}, vol.~93, no.~26, pp.~429--457, 1946.

\bibitem{gellmann}
M.~Gell-Mann, ``A schematic model of baryons and mesons,'' \emph{Physics Letters}, vol.~8, no.~3, pp.~214--215, 1964.

\bibitem{gibbard1973}
A.~Gibbard,
``Manipulation of Voting Schemes: A General Result,''
\emph{Econometrica}, vol.~41, no.~4, pp.~587--601, 1973.

\bibitem{godel} G\"odel, K. (1931). \"Uber formal unentscheidbare S\"atze der Principia Mathematica und verwandter Systeme~I. \emph{Monatshefte f\"ur Mathematik und Physik}, 38, 173--198.

\bibitem{goodstein} Goodstein, R. L. (1944). On the restricted ordinal theorem. \emph{J.\ Symbolic Logic}, 9(2), 33--41.

\bibitem{goodstein1947}
R.~L.~Goodstein,
``Transfinite Ordinals in Recursive Number Theory,''
\emph{Journal of Symbolic Logic} 12(4), 123--129 (1947).

\bibitem{graner1994} F.~Graner and B.~Dubrulle, ``Titius-Bode laws in the solar system. I. Scale invariance explains everything,'' \emph{Astron. Astrophys.}\ 282, 262--268 (1994).

\bibitem{greensite}
J.~Greensite, \emph{An Introduction to the Confinement Problem}, 2nd ed., Lecture Notes in Physics, vol.~972, Springer, 2020. (Center symmetry and $Z(N)$ vortex confinement: Ch.~4--5.)

\bibitem{grothendieck} Grothendieck, A. (1984). \emph{Esquisse d'un Programme}.

\bibitem{grothendieck-rs} Grothendieck, A. (1985--1986). \emph{R\'ecoltes et Semailles: R\'eflexions et T\'emoignage sur un Pass\'e de Math\'ematicien}. Universit\'e des Sciences et Techniques du Languedoc, Montpellier.

\bibitem{grothendieck_ega}
A.~Grothendieck and J.~Dieudonn\'e,
``\'El\'ements de g\'eom\'etrie alg\'ebrique,''
\emph{Inst. Hautes \'Etudes Sci. Publ. Math.}, 1960--1967.

\bibitem{gutin} Gutin, R. (2020). Constructive proof of Herschfeld's Convergence Theorem. \emph{Mathematics of Computation}, 89, 237--251.

\bibitem{hameroff2014} S.~Hameroff and R.~Penrose, ``Consciousness in the universe: A review of the `Orch OR' theory,'' \emph{Physics of Life Reviews}\ 11(1), 39--78 (2014).

\bibitem{hatcher} Hatcher, A. (2002). \emph{Algebraic Topology}. Cambridge University Press.

\bibitem{hod1998}
S.~Hod (1998).
``Bohr's correspondence principle and the area spectrum of quantum black holes.''
\textit{Phys.\ Rev.\ Lett.} 81, 4293.

\bibitem{hopfield} Hopfield, J. J. (1982). Neural networks and physical systems with emergent collective computational abilities. \emph{Proc.\ Natl.\ Acad.\ Sci.}, 79(8), 2554--2558.

\bibitem{hull2018}
J.~C.~Hull,
\emph{Options, Futures, and Other Derivatives},
10th ed., Pearson, 2018.

\bibitem{ilinski1997}
K.~Ilinski,
``Physics of Finance,''
\emph{arXiv:hep-th/9710148}, 1997.

\bibitem{jaffe2000}
A.~Jaffe and E.~Witten,
``Quantum Yang-Mills Theory,''
Official Problem Description, Clay Mathematics Institute (2000).

\bibitem{jaramillo} Jaramillo Salazar, J. D. (2020). Hyperoperations in Exponential Fields. \emph{Journal of Algebra and Its Applications}, 19(1).

\bibitem{jiang_wilczek2019}
Q.-D.~Jiang and F.~Wilczek,
``Chiral Casimir forces: Repulsive, enhanced, tunable,''
\emph{Physical Review B}, vol.~99, no.~12, p.~125403, 2019.

\bibitem{katrin}
M.~Aker \emph{et al.} (KATRIN Collaboration), ``Direct neutrino-mass measurement with sub-electronvolt sensitivity,'' \emph{Nature Physics}, vol.~18, pp.~160--166, 2022.

\bibitem{kerodon} Lurie, J. (2018--). \emph{Kerodon}. \url{https://kerodon.net}.

\bibitem{kim2001}
J.~Kim, ``Neutrino Mass and Mixing,'' \emph{Phys. Rep.} \textbf{150}, 1--177, 2001.

\bibitem{knuth-surreal} Knuth, D. E. (1974). \emph{Surreal Numbers}. Addison-Wesley.

\bibitem{knuth1976}
D.~E.~Knuth,
``Mathematics and Computer Science: Coping with Finiteness,''
\emph{Science} 194(4271), 1235--1242 (1976).
Introduces the up-arrow notation for Goodstein's hyperoperation hierarchy.

\bibitem{kolmogorov1941}
A.~N.~Kolmogorov,
``The local structure of turbulence in incompressible viscous fluid for very large Reynolds numbers,''
\emph{Doklady Akademii Nauk SSSR}, vol.~30, pp.~299--303, 1941.

\bibitem{konkoly2021} K.~R.~Konkoly et al., ``Real-time dialogue between experimenters and dreamers during REM sleep,'' \emph{Current Biology}\ 31(7), 1417-1427.e6 (2021).

\bibitem{krapivin2025optimal}
M.~Farach-Colton, A.~Krapivin, and W.~Kuszmaul,
``Optimal Bounds for Open Addressing Without Reordering,''
arXiv preprint arXiv:2501.02305 (2025).

\bibitem{landauer} Landauer, R. (1961). Irreversibility and Heat Generation in the Computing Process. \emph{IBM Journal of Research and Development}, 5(3), 183--191.

\bibitem{layden2025}
D.~Layden et al., ``Quantum Error Correction Below the Surface Code Threshold,'' \emph{Nature} \textbf{638}, 2025.

\bibitem{lemouël2025}
J.-L.~Le Mou\"{e}l, V.~Courtillot, and S.~Gibert, ``SolDynamo: Solar Dynamo and Planetary Coupling,'' preprint, 2025.

\bibitem{lev-finite-math}
Lev, F.~M. (2020).
Finite Mathematics as the Foundation of Classical Mathematics and Quantum Theory.
Springer.

\bibitem{lev-gfqt}
Lev, F.~M. (2010).
Introduction to A Quantum Theory over A Galois Field.
\emph{Symmetry}, 2(4), 1810--1845.

\bibitem{lev-gfqt-3}
F.~M.~Lev,
``Introduction to A Quantum Theory over A Galois Field,''
\emph{Symmetry} 2(4), 1810--1845 (2010).

\bibitem{lockwood} Lockwood, J. (2024). Chronometric Coordination in Distributed Temporal Systems. \emph{Journal of Distributed Computing}, 37(4), 215--234.

\bibitem{lockwood_zpe2025}
J.~Lockwood,
``Fractal Spacetime and Zero-Point Energy,'' 2025.

\bibitem{lopezdeprado2018}
M.~L\'{o}pez de Prado,
\emph{Advances in Financial Machine Learning},
Wiley, 2018.
Chapter 11: ``The Dangers of Backtesting.'' Sustained precision above 50\% on crash classification tasks invites scrutiny for look-ahead bias and data leakage.

\bibitem{luminet2003}
J.-P.~Luminet, J.~R.~Weeks, A.~Riazuelo, R.~Lehoucq, J.-P.~Uzan (2003).
``Dodecahedral space topology as an explanation for weak wide-angle temperature correlations in the cosmic microwave background.''
\textit{Nature} 425, 593--595.

\bibitem{lurie-htt} Lurie, J. (2009). \emph{Higher Topos Theory}. Princeton University Press.

\bibitem{maclane-moerdijk} Mac Lane, S. \& Moerdijk, I. (1992). \emph{Sheaves in Geometry and Logic}. Springer-Verlag.

\bibitem{malcev} Mal'cev, A. I. (1955). Analytic loops. \emph{Mat.\ Sb.}, 36(78), 569--576.

\bibitem{maluf2013}
J.~Maluf, ``The Teleparallel Equivalent of General Relativity,'' \emph{Annalen der Physik} \textbf{525}, 339--357, 2013.

\bibitem{mandelbrot} Mandelbrot, B. B. (1982). \emph{The Fractal Geometry of Nature}. W.\,H.\ Freeman.

\bibitem{mandelbrot1963}
B.~B.~Mandelbrot,
``The Variation of Certain Speculative Prices,''
\emph{The Journal of Business}, vol.~36, no.~4, pp.~394--419, 1963.

\bibitem{manton2002}
N.~S.~Manton,
``Fluid mechanical models of financial markets,''
Working paper, DAMTP Cambridge, 2002.

\bibitem{massot-blueprints} Massot, P. (2023). Lean Blueprints: Bridging Informal and Formal Mathematics. In \emph{ITP 2023}, LIPIcs.

\bibitem{mathlib} The mathlib Community. (2020). The Lean Mathematical Library. In \emph{CPP '20}, pp.~367--381. ACM.

\bibitem{may-ponto} May, J. P. \& Ponto, K. (2012). \emph{More Concise Algebraic Topology}. University of Chicago Press.

\bibitem{mckinney05} J.~McKinney et al., ``A revised mechanism for tryptophan hydroxylase,'' \emph{J.~Biol.~Chem.}\ 280, 13265 (2005).

\bibitem{milne-rhff}
Milne, J.~S. (2015).
The Riemann Hypothesis over Finite Fields: From Weil to the Present Day (a conjecture resolved for function fields by Deligne).
In \emph{The Legacy of Bernhard Riemann After One Hundred and Fifty Years}. International Press.

\bibitem{milne-rhff-12} J.~S.~Milne, ``The Riemann Hypothesis over Finite Fields,'' in \emph{The Legacy of Bernhard Riemann After One Hundred and Fifty Years}, International Press (2015).

\bibitem{milne-rhff-3}
J.~S.~Milne,
``The Riemann Hypothesis over Finite Fields: From Weil to the Present Day,''
in \emph{The Legacy of Bernhard Riemann After One Hundred and Fifty Years},
International Press (2015).

\bibitem{milne_cft}
J.~S.~Milne,
``Class Field Theory,''
lecture notes, v4.03, 2020.

\bibitem{milnor} Milnor, J. (1963). \emph{Morse Theory}. Princeton University Press.

\bibitem{montgomery} Montgomery, H. L. (1973). The pair correlation of zeros of the zeta function. \emph{Analytic Number Theory}, 181--193.

\bibitem{motl2003}
L.~Motl and A.~Neitzke (2003).
``Asymptotic black hole quasinormal frequencies.''
\textit{Adv.\ Theor.\ Math.\ Phys.} 7, 307--330.

\bibitem{mrs_glutamate} R.~A.~E.~Edden et al., ``Proton MR spectroscopy of GABA, glutamate, and glutamine,'' \emph{Neuroimage}, 2012.

\bibitem{mrs_variance} A.~Bogaschewsky et al., ``Reproducibility of Glx and GABA measurements in the human brain,'' \emph{Magnetic Resonance in Medicine}, 2019.

\bibitem{nagatsu64} T.~Nagatsu, M.~Levitt, and S.~Udenfriend, ``Tyrosine Hydroxylase: The Initial Step in Norepinephrine Biosynthesis,'' \emph{J Biol Chem}\ 239, 2910--2917 (1964).

\bibitem{nash} Nash, J. F. (1950). Equilibrium points in $n$-person games. \emph{Proc.\ Natl.\ Acad.\ Sci.}, 36(1), 48--49.

\bibitem{neukirch}
J.~Neukirch,
\emph{Algebraic Number Theory},
Grundlehren der mathematischen Wissenschaften, vol.~322, Springer, 1999.

\bibitem{newton}
I.~Newton, \emph{Opticks: or, A Treatise of the Reflexions, Refractions, Inflexions and Colours of Light}, 1704.

\bibitem{odlyzko} Odlyzko, A. M. (1987). On the distribution of spacings between zeros of the zeta function. \emph{Mathematics of Computation}, 48(177), 273--308.

\bibitem{odrzywolek} Odrzywolek, A. (2026). All elementary functions from a single operator. \emph{arXiv preprint}.

\bibitem{padicCognition} A.~Khrennikov, \emph{Information Dynamics in Cognitive, Psychological, Social and Anomalous Phenomena}, Springer (2004).

\bibitem{penrose} Penrose, R. (1989). \emph{The Emperor's New Mind}. Oxford University Press.

\bibitem{penrose-qp} Penrose, R. (1974). The role of aesthetics in pure and applied mathematical research. \emph{Bull.\ IMA}, 10(7/8), 266--271.

\bibitem{perlin1985image}
Ken Perlin,
\textit{An Image Synthesizer},
SIGGRAPH Computer Graphics, 19(3):287--296, 1985.

\bibitem{persson_minecraft}
M. Persson et al.,
\textit{Minecraft: Deterministic Voxel Generation and LCGs},
Mojang Specifications (Informal Technical Documentation), 2009--2011.

\bibitem{pitkanenTGD} M.~Pitk\"anen, ``TGD Inspired Theory of Consciousness,'' \emph{Journal of Non-Locality and Remote Mental Interactions}\ 1 (2002).

\bibitem{planck2020}
Planck Collaboration (2020).
``Planck 2018 results.\ VII.\ Isotropy and statistics of the CMB.''
\textit{Astron.\ Astrophys.} 641, A7.

\bibitem{pollard1971}
J.~M.~Pollard,
``The fast Fourier transform in a finite field,''
\emph{Math.\ Comp.} 25(114), 365--374 (1971).

\bibitem{ptvv} Pantev, T., To\"en, B., Vaqui\'e, M. and Vezzosi, G. (2013). Shifted symplectic structures. \emph{Publ. Math. Inst. Hautes \'Etudes Sci.} 117, 271--328.

\bibitem{rado1962}
T.~Rad\'o,
``On Non-Computable Functions,''
\emph{Bell System Technical Journal} 41(3), 877--884 (1962).

\bibitem{riemann} Riemann, B. (1859). Ueber die Anzahl der Primzahlen unter einer gegebenen Gr\"osse. \emph{Monatsberichte der Berliner Akademie}.

\bibitem{robinson} Robinson, A. (1966). \emph{Non-standard Analysis}. North-Holland.

\bibitem{roy1954} A.~E.~Roy and M.~W.~Ovenden, ``On the occurrence of commensurable mean motions in the solar system,'' \emph{MNRAS}\ 114, 232--241 (1954).

\bibitem{safe-verification} Xin, H. et al. (2025). Safe: Step-aware formal verification for LLM reasoning. \emph{Proc.\ ACL}, pp.~4821--4838.

\bibitem{satterthwaite1975}
M.~A.~Satterthwaite,
``Strategy-Proofness and Arrow's Conditions,''
\emph{Journal of Economic Theory}, vol.~10, no.~2, pp.~187--217, 1975.

\bibitem{savitz20} J.~Savitz, ``The kynurenine pathway: a finger in every pie,'' \emph{Mol Psychiatry}\ 25, 131--147 (2020).

\bibitem{schafer} Schafer, R. D. (1966). \emph{An Introduction to Nonassociative Algebras}. Academic Press.

\bibitem{schrodinger}
E.~Schr\"odinger, ``Quantisierung als Eigenwertproblem,'' \emph{Annalen der Physik}, vol.~384, no.~4, pp.~361--376, 1926.

\bibitem{schrodinger44}
E.~Schr\"odinger,
\emph{What is Life? The Physical Aspect of the Living Cell},
Cambridge University Press, 1944.
Chapter~6: ``Order, Disorder, and Entropy.''

\bibitem{schwieler16} L.~Schwieler et al., ``Electroconvulsive therapy suppresses the neurotoxic branch of the kynurenine pathway in treatment-resistant depressed patients,'' \emph{Transl Psychiatry}\ 6, e906 (2016).

\bibitem{serre_local}
J.-P.~Serre,
\emph{Local Fields},
Graduate Texts in Mathematics, vol.~67, Springer, 1979.

\bibitem{shannon} Shannon, C. E. (1948). A mathematical theory of communication. \emph{Bell System Technical Journal}, 27(3), 379--423.

\bibitem{shannon48}
C.~E.~Shannon,
``A Mathematical Theory of Communication,''
\emph{Bell System Technical Journal} 27, 379--423, 623--656 (1948).

\bibitem{shor1994}
P.~W.~Shor,
``Algorithms for quantum computation: discrete logarithms and factoring,''
\emph{Proceedings 35th Annual Symposium on Foundations of Computer Science}, IEEE, 124--134 (1994).

\bibitem{shulgin_pihkal}
A.~T.~Shulgin and A.~Shulgin,
\emph{PiHKAL: A Chemical Love Story},
Transform Press, 1991.

\bibitem{sono96} M.~Sono et al., ``Heme-containing oxygenases,'' \emph{Chem.\ Rev.}\ 96, 2841 (1996).

\bibitem{srinivasan15} B.~Srinivasan et al., ``TEER measurement techniques for in vitro barrier model systems,'' \emph{JALA}\ 20, 107 (2015).

\bibitem{steenrod} Steenrod, N. (1951). \emph{The Topology of Fibre Bundles}. Princeton University Press.

\bibitem{stillinger1977}
F.~Stillinger, ``Axiomatic Basis for Spaces with Noninteger Dimension,'' \emph{J. Math. Phys.} \textbf{18}, 1224--1234, 1977.

\bibitem{stockigt11} J.~St\"ockigt et al., ``The Pictet--Spengler reaction in nature and in organic chemistry,'' \emph{Angew.~Chem.}\ 50, 7580 (2011).

\bibitem{takahashi2019}
S.~Takahashi, Y.~Chen, and K.~Tanaka-Ishii,
``Modeling Financial Time-Series with Generative Adversarial Networks,''
\emph{Physica A: Statistical Mechanics and its Applications}, vol.~527, 121261, 2019.

\bibitem{tarski} Tarski, A. (1936). Der Wahrheitsbegriff in den formalisierten Sprachen. \emph{Studia Philosophica}, 1, 261--405.

\bibitem{thooft1981}
G.~'t~Hooft,
``Topology of the gauge condition and new confinement phases in non-Abelian gauge theories,''
\emph{Nucl. Phys. B}\ 190, 455--478 (1981).

\bibitem{timegan2019}
J.~Yoon, D.~Jarrett, and M.~van der Schaar,
``Time-series Generative Adversarial Networks,''
\emph{Advances in Neural Information Processing Systems (NeurIPS)}, 2019.

\bibitem{tsai-spaceq} Eby, J., Safronova, M. S., \& Tsai, Y.-D. (2022). Direct detection of ultralight dark matter bound to the Sun with space quantum sensors. \emph{Nature Astronomy}, 6, 1233--1239.

\bibitem{ttsgan2024}
X.~Li, V.~Metsis, H.~Wang, and A.~Ngu,
``TTS-GAN: A Transformer-based Time-Series Generative Adversarial Network,''
\emph{arXiv:2202.02691v3}, 2024.

\bibitem{varieschi2020}
G.~Varieschi, ``Newtonian Fractional-Dimension Gravity and MOND,'' \emph{Found. Phys.} \textbf{50}, 1608--1644, 2020.

\bibitem{veblen} Veblen, O. (1908). Continuous Increasing Functions of Finite and Transfinite Ordinals. \emph{Trans.\ AMS}, 9(3), 280--292.

\bibitem{wiener} Wiener, N. (1948). \emph{Cybernetics: Or Control and Communication in the Animal and the Machine}. MIT Press.

\bibitem{wilczek1973}
F.~Wilczek, ``Ultraviolet Stability of the Vacuum in a Non-Abelian Gauge Theory,'' \emph{Phys. Rev. Lett.} \textbf{30}, 1343--1346, 1973.

\bibitem{wilson1974}
K.~G.~Wilson,
``Confinement of quarks,''
\emph{Phys. Rev. D}\ 10, 2445--2459 (1974).


% ────────────────────────────────────────────
% AUTHOR'S WORKS, PREPRINTS, AND
% PRIVATE COMMUNICATIONS
% ────────────────────────────────────────────
% Works by the author(s) of this volume,
% preprints not yet peer-reviewed, private
% communications, and informal sources.
% These have NOT undergone external peer review.
% ────────────────────────────────────────────

\bibitem{Aingel}
La Rue, D. (2026).
\texttt{Aingel\_proofs}: The ZKP Holochain Envelope over the Metareal Manifold.
Formal algebraic synthesis.

\bibitem{GoldenTetrahedron}
La Rue, D. (2026).
\textit{The Golden Tetrahedron: Gauge Center Hierarchy and Chronometric Structure}.

\bibitem{InfochemicalTime}
La Rue, D. (2026).
\texttt{InfochemicalTime\_proofs}: Quasi-Epi-Periodic Geometric Time.
Formal algebraic synthesis.

\bibitem{MetarealManifold}
La Rue, D. (2026).
\texttt{MetarealManifold\_proofs}: The 12-Dimensional Observer-Manifold Space.
Formal algebraic synthesis.

\bibitem{PrimeFieldMechanics}
La Rue, D. (2026).
\texttt{PrimeFieldMechanics\_proofs}: Agentic Mechanics.
Formal algebraic synthesis.

\bibitem{bionetic_ambrosia}
D.~La Rue,
``The Bionetic Ambrosia Protocol,''
June 2026.

\bibitem{btc_indicator2025}
Various,
``Bitcoin as a Leading Indicator for Risk Appetite: 24/7 Liquidity Signals and Equity Market Correlation,''
Market microstructure research and institutional reports, 2024--2026.

\bibitem{crashbenchmark2024}
Various,
``Stock Market Crash Prediction: Evaluation Metrics and Model Comparison,''
Academic surveys across machine learning and econometric approaches, 2020--2024.
Typical precision for 3-month crash prediction: $\sim$15\% at recall $\sim$59\%.

\bibitem{dwm}
D.~La Rue,
``Digital Wave Mechanics,'' 2025.

\bibitem{dwm_fst}
J.~Lockwood and D.~La Rue, ``Digital Wave Mechanics and Fractal Spacetime: Adelic Orbit Structure of the Chronometric Triangle,'' preprint, June 2026.

\bibitem{ewsbenchmark2024}
Various,
``Machine Learning Early Warning Systems for Financial Crises,''
Studies from Bank of Canada, ECB, IMF, and academic papers, 2020--2026.
State-of-the-art AUROC: 0.75--0.88 for 6-month recession forecasting using 15--50 macroeconomic indicators.

\bibitem{gans_diffusion2024}
Various,
``GAN-Guided Diffusion Models for Financial Time Series,''
\emph{arXiv preprints}, 2024--2025.

\bibitem{golden_tetrahedron}
D.~La Rue,
``The Golden Tetrahedron: Algebraic Foundations of Standard Model Symmetries,''
2026.

\bibitem{hci_prompt2025}
Various,
``Human-Computer Interaction in Financial Trading Systems: Market Research and Design Principles,''
Industry reports and academic surveys, 2025--2026.

\bibitem{larue} La Rue, D. (2026). Protoreal Zeta. Formal verification of topological and algebraic axioms. \url{https://github.com/Dielawn-01/ProtorealZeta}.

\bibitem{larue-cde} La Rue, D. \& Lockwood, J. (2026). SU(3) Center Drift Extension: Golden Cosets and Temporal Spectral Dynamics.

\bibitem{larue2026chromatic}
D.~La Rue.
\textit{The SU(3) Center Triangle and Fractal SpaceTime} (Digital Wave Mechanics).
Spectrality Institute, 2026.

\bibitem{larue2026lc}
D. La Rue.
\textit{Logical Creativity and the La Rue Manifold}.
Spectrality Institute, 2026.

\bibitem{larue_chromatic}
D.~La Rue, ``The SU(3) Center Triangle: Golden Split Primes, Standing Waves in Digit Space, and the Dark-Bright Palindrome Resonance,'' preprint, June 2026.

\bibitem{larue_dwm}
D.~La Rue and J.~Lockwood,
``Digital Wave Mechanics and Fractal Spacetime,''
Preprint, June 2026.

\bibitem{larue_lc}
D.~La Rue, ``Logical Creativity: A Discussion on the Foundations of Mathematics,'' preprint, June 2026.

\bibitem{lc} D.~La Rue, \emph{Logical Creativity}, Protoreal Zeta Series, 2026.

\bibitem{lockwood2025dwm}
D.~La Rue and J.~Lockwood.
\textit{Digital Wave Mechanics}.
Spectrality Institute, 2025.

\bibitem{lockwood2025zpe}
J.~Lockwood.
\textit{Zero-Point Energy: First Principles, Units, and Observables}.
Spectrality Institute, Sept 2025.

\bibitem{lockwood_chrono}
J.~Lockwood,
``Chronometric Constants and Frozen Spectral Relations,''
Private communication, 2026.

\bibitem{lockwood_chrono3}
J.~Lockwood, ``Chronometric Constants and Frozen Spectral Relations (Chrono~3),'' private communication, 2026.

\bibitem{lockwood_chrono4}
J.~Lockwood, ``Chronometric Constants and Frozen Spectral Relations (Chrono~4),'' private communication, 2026.

\bibitem{lockwood_larue_dwm}
J.~Lockwood and D.~La Rue,
``Digital Wave Mechanics: From Newton's Prism to Standing Waves
in Finite Golden Fields,'' 2026.

\bibitem{lockwood_v11}
M.~Lockwood, ``Solar Cycle Prediction: v11 Update,'' private communication, 2025.

\bibitem{metareal_asi}
D.~La Rue,
``Metareal ASI Chromodynamics: The Golden Veblen Resolution of G\"odelian Incompleteness,''
June 2026.

\bibitem{plasma} D.~La Rue, \emph{InfoPhys Lattice Dynamics}, Metamaterial Research Export, 2026. \url{https://github.com/Dielawn-01/InfoPhys}

\bibitem{protoreal} D.~La Rue, \emph{Protoreal Zeta}, formal axioms, 2026. \url{https://github.com/Dielawn-01/Protoreal_Zeta}

\bibitem{rj2} J.~Lockwood and D.~Hansley, \emph{Royal Jelly: Field--Transport Foundations for Bioelectromagnetics}, Spectrality Institute, 2025.

\bibitem{rja} D.~La Rue, \emph{Transport--Algebraic Correspondence in Bioelectromagnetics: An Analysis of the Royal Jelly Framework}, InfoPhys Axioms, 2026.

\bibitem{steiniger2026epg}
Steiniger, R.,
``Engineering Persistent Geometric Identities in LLMs,''
Preprint v3, May 2026.

\bibitem{walker_infoton}
M.~Walker, ``The Infoton, the Blackbody Constant of Information-Energy, and the
Landauer Wall,'' private communication, 2024.

% ═══ Time Crystals ═══════════════════════════════════
\bibitem{wilczek2012}
F.~Wilczek,
``Quantum Time Crystals,''
\emph{Physical Review Letters}, vol.~109, 160401, 2012.

\bibitem{watanabe2015}
H.~Watanabe and M.~Oshikawa,
``Absence of Quantum Time Crystals,''
\emph{Physical Review Letters}, vol.~114, 251603, 2015.

\bibitem{zhang2017}
J.~Zhang et al.,
``Observation of a Discrete Time Crystal,''
\emph{Nature}, vol.~543, pp.~217--220, 2017.

\bibitem{mi2022}
X.~Mi et al.,
``Time-crystalline eigenstate order on a quantum processor,''
\emph{Nature}, vol.~601, pp.~531--536, 2022.

\bibitem{else2016}
D.~V.~Else, B.~Bauer, and C.~Nayak,
``Floquet Time Crystals,''
\emph{Physical Review Letters}, vol.~117, 090402, 2016.

% ═══ Przybylski's Star ═══════════════════════════════
\bibitem{cowley2004}
C.~R.~Cowley, W.~P.~Bidelman, S.~Hubrig, G.~Mathys, and D.~J.~Bord,
``On the Possible Presence of Promethium in the Spectrum of HD~101065 (Przybylski's Star),''
\emph{The Astrophysical Journal}, vol.~603, pp.~354--367, 2004.

\bibitem{gopka2004}
V.~F.~Gopka, A.~V.~Yushchenko, A.~V.~Shavrina et al.,
``Identification of absorption lines of the transuranium elements in the spectrum of Przybylski's star (HD~101065),''
\emph{Kinematika i Fizika Nebesnykh Tel}, vol.~20, pp.~210--220, 2004.

\bibitem{delbaen1994}
F.~Delbaen and W.~Schachermayer,
``A general version of the fundamental theorem of asset pricing,''
\emph{Mathematische Annalen}, vol.~300, pp.~463--520, 1994.

\bibitem{wyckoff1931}
R.~D.~Wyckoff,
\emph{The Richard D.~Wyckoff Method of Trading and Investing in Stocks}.
New York: Wyckoff Associates, 1931.

\bibitem{weil1948courbes}
A.~Weil,
\emph{Sur les courbes alg\'ebriques et les vari\'et\'es qui s'en d\'eduisent}.
Paris: Hermann, 1948.

% ═══ Algorithmic Trading / FTAP ════════════════════════
\bibitem{harrison_kreps1979}
J.~M.~Harrison and D.~M.~Kreps,
``Martingales and arbitrage in multiperiod securities markets,''
\emph{Journal of Economic Theory}, vol.~20, no.~3, pp.~381--408, 1979.

\bibitem{harrison_pliska1981}
J.~M.~Harrison and S.~R.~Pliska,
``Martingales and stochastic integrals in the theory of continuous trading,''
\emph{Stochastic Processes and their Applications}, vol.~11, no.~3, pp.~215--260, 1981.

\bibitem{doob1953}
J.~L.~Doob,
\emph{Stochastic Processes},
Wiley, 1953.
Chapter~VII: Martingale theory and the tower property of conditional expectations.

\bibitem{niederreiter1992}
H.~Niederreiter,
\emph{Random Number Generation and Quasi-Monte Carlo Methods},
CBMS-NSF Regional Conference Series in Applied Mathematics, vol.~63, SIAM, 1992.
Chapters~8--9: inversive congruential generators and their equidistribution properties.

\bibitem{avellaneda_lee2010}
M.~Avellaneda and J.-H.~Lee,
``Statistical Arbitrage in the U.S.\ Equities Market,''
\emph{Quantitative Finance}, vol.~10, no.~7, pp.~761--782, 2010.

\bibitem{hurst_ooi_pedersen2017}
B.~Hurst, Y.~H.~Ooi, and L.~H.~Pedersen,
``A Century of Evidence on Trend-Following Investing,''
\emph{Journal of Portfolio Management}, vol.~44, no.~1, pp.~15--29, 2017.

\bibitem{qian2005}
E.~E.~Qian,
``Risk Parity Portfolios: Efficient Portfolios Through True Diversification,''
PanAgora Asset Management, White Paper, 2005.

\bibitem{embrechts2002}
P.~Embrechts, A.~McNeil, and D.~Straumann,
``Correlation and Dependence in Risk Management: Properties and Pitfalls,''
in \emph{Risk Management: Value at Risk and Beyond}, ed.\ M.~Dempster,
Cambridge University Press, pp.~176--223, 2002.

% ────────────────────────────────────────────────────────────────
% TEMPORAL QUASICRYSTALS & BIOLOGICAL NON-ASSOCIATIVITY
% ────────────────────────────────────────────────────────────────

\bibitem{penrose1974}
R.~Penrose,
``The Role of Aesthetics in Pure and Applied Mathematical Research,''
\emph{Bulletin of the Institute of Mathematics and its Applications},
vol.~10, no.~7/8, pp.~266--271, 1974.
\textit{Introduces aperiodic tilings with golden-ratio structure (Penrose tilings),
the spatial precursor to quasicrystals.}

\bibitem{schrodinger1944}
E.~Schr\"odinger,
\emph{What is Life? The Physical Aspect of the Living Cell},
Cambridge University Press, 1944.
\textit{Predicts that the genetic material is an ``aperiodic crystal''---
a structure with long-range order but no periodic repetition.}

\bibitem{shechtman1984}
D.~Shechtman, I.~Blech, D.~Gratias, and J.~W.~Cahn,
``Metallic Phase with Long-Range Orientational Order and No Translational Symmetry,''
\emph{Physical Review Letters}, vol.~53, no.~20, pp.~1951--1953, 1984.
\textit{Discovery of quasicrystals. Nobel Prize in Chemistry, 2011.}

\bibitem{crick1966}
F.~H.~C.~Crick,
``Codon--Anticodon Pairing: The Wobble Hypothesis,''
\emph{Journal of Molecular Biology}, vol.~19, no.~2, pp.~548--555, 1966.
\textit{The wobble position allows non-canonical base pairing at the third codon
position, introducing controlled imperfection into the genetic code.}

\bibitem{yao2017}
N.~Y.~Yao, A.~C.~Potter, I.-D.~Potirniche, and A.~Vishwanath,
``Discrete Time Crystals: Rigidity, Criticality, and Realizations,''
\emph{Physical Review Letters}, vol.~118, no.~3, p.~030401, 2017.
\textit{Theoretical blueprint for discrete time crystals, defining conditions
for robust subharmonic order in Floquet systems.}

% ────────────────────────────────────────────────────────────────
% GAUSSIAN RELATIVITY / ADELIC FORCE LAW
% ────────────────────────────────────────────────────────────────

\bibitem{lemouël2025ext}
J.-L.~Le Mou\"{e}l, F.~Lopes, V.~Courtillot, and D.~Gibert,
``Planetary Forcing of Solar Activity: A Review,''
preprint (SolDynamoExtForcing), 2025.
\textit{Identifies 11 shared pseudo-periods between sunspot number (SSN, 1750--2025)
and length-of-day (LOD, 1780--2025) via Singular Spectrum Analysis.
Correlation $r > 0.999$, bootstrap $p \leq 10^{-3}$.}

\bibitem{holme2013}
R.~Holme and O.~de~Viron,
``Characterization and implications of intradecadal variations in length of day,''
\emph{Nature}, vol.~499, pp.~202--204, 2013.
\textit{LOD analysis establishing electromagnetic core-mantle coupling coefficient
$\approx 0.80$ for Earth.}

\bibitem{dziewonski1981}
A.~M.~Dziewonski and D.~L.~Anderson,
``Preliminary reference Earth model,''
\emph{Physics of the Earth and Planetary Interiors}, vol.~25, no.~4, pp.~297--356, 1981.
\textit{PREM. Core radius $R_{\text{core}} = 3480$~km, planet radius $R_\oplus = 6371$~km,
ratio $= 0.546$.}

\bibitem{pettengill1965}
G.~H.~Pettengill and R.~B.~Dyce,
``A radar determination of the rotation of the planet Mercury,''
\emph{Nature}, vol.~206, pp.~1240, 1965.
\textit{Discovery of Mercury's 3:2 spin-orbit resonance. Tidal lock zeroes the
conjugate golden channel $\bar{\varphi} = (\omega - \iota)/2 = 0$.}

\bibitem{alken2021}
P.~Alken et al.,
``International Geomagnetic Reference Field: the thirteenth generation,''
\emph{Earth, Planets and Space}, vol.~73, art.~49, 2021.
\textit{IGRF-13. Earth dipole tilt $= 11.5^\circ \pm 0.3^\circ$.}

\bibitem{connerney2022}
J.~E.~P.~Connerney et al.,
``A new model of Jupiter's magnetic field at the completion of Juno's prime mission,''
\emph{Journal of Geophysical Research: Planets}, vol.~127, e2021JE007055, 2022.
\textit{JRM33 model. Jupiter dipole tilt $= 10.3^\circ \pm 0.5^\circ$.}

\bibitem{dougherty2018}
M.~K.~Dougherty et al.,
``Saturn's magnetic field revealed by the Cassini Grand Finale,''
\emph{Science}, vol.~362, eaat5434, 2018.
\textit{Cassini MAG. Saturn dipole tilt $< 0.01^\circ$ (axisymmetric),
consistent with hexagonal ($k = 6$) symmetry lock.}

\bibitem{ness1986}
N.~F.~Ness, M.~H.~Acu\~na, K.~W.~Behannon, L.~F.~Burlaga,
J.~E.~P.~Connerney, and R.~P.~Lepping,
``Magnetic fields at Uranus,''
\emph{Science}, vol.~233, pp.~85--89, 1986.
\textit{Voyager~2. Uranus dipole tilt $58.6^\circ \pm 2^\circ$ from rotation axis,
offset 0.3 planetary radii from center.}

\bibitem{ness1989}
N.~F.~Ness, M.~H.~Acu\~na, L.~F.~Burlaga, J.~E.~P.~Connerney,
and R.~P.~Lepping,
``Magnetic fields at Neptune,''
\emph{Science}, vol.~246, pp.~1473--1478, 1989.
\textit{Voyager~2. Neptune dipole tilt $47^\circ \pm 3^\circ$ from rotation axis,
offset 0.55 planetary radii from center.}

\bibitem{kivelson2002}
M.~G.~Kivelson, K.~K.~Khurana, and M.~Volwerk,
``The permanent and inductive magnetic moments of Ganymede,''
\emph{Icarus}, vol.~157, pp.~507--522, 2002.
\textit{Galileo. Ganymede intrinsic dipole anti-aligned with Jupiter's field
($\sim 176^\circ$). Only moon with self-sustaining magnetosphere.}

\bibitem{kramers40}
H.~A.~Kramers,
``Brownian motion in a field of force and the diffusion model of chemical reactions,''
\emph{Physica}, vol.~7, no.~4, pp.~284--304, 1940.
\textit{The foundational derivation of the overdamped Kramers escape rate for barrier crossing in the presence of thermal noise.}

\bibitem{snyder65}
S.~H.~Snyder and C.~R.~Merrill,
``A relationship between the hallucinogenic activity of drugs and their electronic configuration,''
\emph{Proc.\ Natl.\ Acad.\ Sci.\ USA}, vol.~54, no.~1, pp.~258--266, 1965.
\textit{Pioneering QSAR analysis linking psychedelic potency to HOMO energy and superdelocalizability of the indole $\pi$-system.}

\bibitem{FDA_austedo}
U.S.~Food and Drug Administration,
``FDA approves first deuterated drug---Austedo (deutetrabenazine) for Huntington's disease chorea,''
FDA Press Release, April 2017.
\textit{First approved deuterium-modified drug. Six ${}^1$H$\to$${}^2$D substitutions reduce CYP2D6 metabolism via kinetic isotope effect (KIE $\approx 2$--$5$).}

% ────────────────────────────────────────────────────────────────
% ENTRIES ADDED BY 4C AUDIT (2026-07-03)
% ────────────────────────────────────────────────────────────────

\bibitem{albert}
A.~A.~Albert,
``Power-Associative Rings,''
\emph{Trans.\ Amer.\ Math.\ Soc.}, vol.~64, no.~3, pp.~552--593, 1948.
\textit{Foundational classification of power-associative algebras. Establishes that Jordan algebras and alternative algebras are power-associative.}

\bibitem{assem-simson-skowronski}
I.~Assem, D.~Simson, and A.~Skowro\'nski,
\emph{Elements of the Representation Theory of Associative Algebras},
London Mathematical Society Student Texts, vol.~65, Cambridge University Press, 2006.
\textit{Standard reference for quiver representations and Auslander--Reiten theory.}

\bibitem{linrado1965}
T.~Lin and C.~T.~Rad\'o,
``On Non-Computable Functions and Computability Theory,''
\emph{Notices of the AMS}, 1965.

\bibitem{newman2006}
M.~E.~J.~Newman,
``Modularity and community structure in networks,''
\emph{Proc.\ Natl.\ Acad.\ Sci.\ USA}, vol.~103, no.~23, pp.~8577--8582, 2006.

\bibitem{wall1960}
C.~T.~C.~Wall,
``Determination of the Cobordism Ring,''
\emph{Annals of Mathematics}, vol.~72, no.~2, pp.~292--311, 1960.
\textit{Classification of oriented cobordism ring. Foundational for topological field theories.}

\bibitem{watts-strogatz1998}
D.~J.~Watts and S.~H.~Strogatz,
``Collective dynamics of `small-world' networks,''
\emph{Nature}, vol.~393, pp.~440--442, 1998.

\bibitem{landauer61}
R.~Landauer,
``Irreversibility and Heat Generation in the Computing Process,''
\emph{IBM Journal of Research and Development}, vol.~5, no.~3, pp.~183--191, 1961.
\textit{Establishes the Landauer limit: $k_B T \ln 2$ per bit erasure.}

\bibitem{beggs2003}
J.~M.~Beggs and D.~Plenz,
``Neuronal Avalanches in Neocortical Circuits,''
\emph{Journal of Neuroscience}, vol.~23, no.~35, pp.~11167--11177, 2003.
\textit{Discovery of neuronal avalanches following power-law distributions, evidence for self-organized criticality in cortex.}

\bibitem{bethe1939}
H.~A.~Bethe,
``Energy Production in Stars,''
\emph{Physical Review}, vol.~55, no.~5, pp.~434--456, 1939.
\textit{CNO cycle and pp-chain. Nobel Prize in Physics 1967.}

\bibitem{blackmond2024}
D.~G.~Blackmond,
``The Origin of Biological Homochirality,''
\emph{Cold Spring Harbor Perspectives in Biology}, vol.~2, a002147, 2010.
\textit{Review of chiral symmetry breaking mechanisms in prebiotic chemistry.}

\bibitem{caspar1962}
D.~L.~D.~Caspar and A.~Klug,
``Physical Principles in the Construction of Regular Viruses,''
\emph{Cold Spring Harbor Symposia on Quantitative Biology}, vol.~27, pp.~1--24, 1962.
\textit{Caspar--Klug theory of icosahedral virus capsid structure. Triangulation numbers.}

\bibitem{cronin2023}
L.~Cronin and S.~I.~Walker,
``Beyond prebiotic chemistry,''
\emph{Science}, vol.~352, no.~6290, pp.~1174--1175, 2016.
\textit{Assembly theory and the quantification of molecular complexity.}

\bibitem{kauffman2024}
S.~A.~Kauffman,
\emph{At Home in the Universe: The Search for Laws of Self-Organization and Complexity},
Oxford University Press, 1995.

\bibitem{lu1991}
P.~J.~Lu and P.~J.~Steinhardt,
``Decagonal and Quasi-Crystalline Tilings in Medieval Islamic Architecture,''
\emph{Science}, vol.~315, no.~5815, pp.~1106--1110, 2007.

\bibitem{qt45_2026}
D.~La Rue,
``Temporal Quasicrystals and the QT-45 Hypothesis,''
preprint, 2026.

\bibitem{tassinari2019}
F.~Tassinari et al.,
``Enhanced Hydrogen Production from a Chiral Conducting Polymer,''
\emph{ACS Nano}, vol.~13, no.~8, pp.~8845--8852, 2019.
\textit{Experimental demonstration of chirality-induced spin selectivity (CISS) effect.}

\bibitem{weinberg_qft}
S.~Weinberg,
\emph{The Quantum Theory of Fields, Volume~I: Foundations},
Cambridge University Press, 1995.
\textit{Derivation of the $7/8$ spin-statistics factor for fermionic vacuum energy density (\S5.3).}

% ────────────────────────────────────────────────────────────────
% ENTRIES ADDED FOR Ch.19 (ASI Chip)
% ────────────────────────────────────────────────────────────────

\bibitem{EngelFleming07}
G.~S.~Engel, T.~R.~Calhoun, E.~L.~Read, T.-K.~Ahn, T.~Man\v{c}al, Y.-C.~Cheng,
R.~E.~Blankenship, and G.~R.~Fleming,
``Evidence for wavelike energy transfer through quantum coherence in photosynthetic systems,''
\emph{Nature}, vol.~446, pp.~782--786, 2007.

\bibitem{KohmotoKadanoffTang83}
M.~Kohmoto, L.~P.~Kadanoff, and C.~Tang,
``Localization problem in one dimension: Mapping and escape,''
\emph{Phys.\ Rev.\ Lett.}, vol.~50, no.~23, pp.~1870--1872, 1983.

\bibitem{Shechtman84}
D.~Shechtman, I.~Blech, D.~Gratias, and J.~W.~Cahn,
``Metallic Phase with Long-Range Orientational Order and No Translational Symmetry,''
\emph{Physical Review Letters}, vol.~53, no.~20, pp.~1951--1953, 1984.
\textit{Discovery of quasicrystals.\ Nobel Prize in Chemistry, 2011.}

\bibitem{lockwood_gmf34}
J.~Lockwood,
``GMF-34: 34-Gradient Plasma Materials-Synthesis Forge,''
schematics v1/v2/v8, private communication, 2026.
\textit{34-gradient radial layout with central synthesis volume, MHD governing equations, and 4-band acoustic transducer architecture.}

\bibitem{RussellDiagram}
D.~La Rue,
\texttt{RussellDiagram.lean}: Russell Octave Harmonic Structure in Lean~4,
InfoPhysAxioms, 2026.
\textit{Formalizes the Russell 12-tone cycle, carbon--silicon umbral shift, and the \texttt{RussellOctaveNode} $\to$ \texttt{ProtorealManifold} projection.}

\bibitem{Mizutani2017}
U.~Mizutani and H.~Sato,
``The Physics of the Hume-Rothery Electron Concentration Rule,''
\emph{Crystals}, vol.~7, no.~1, art.~9, 2017.
\textit{Comprehensive derivation of e/a = 7/4 for icosahedral QCs via Friedel's scattering framework. Establishes the Jones zone / Fermi surface matching for stable QC compositions.}

\bibitem{NegTriangularity2026}
A.~Pochelon et al.,
``Enhanced confinement in negative-triangularity shaped tokamak plasmas,''
\emph{Plasma Physics and Controlled Fusion}, Outstanding Paper Prize, 2026.
\textit{Negative-triangularity plasma shaping for improved MHD stability.}

\bibitem{PlasmacousticMeta2025}
Various,
``Plasmacoustic Metalayers: Corona Discharge Transducers for Broadband Active Acoustic Metamaterials,''
in Proc.\ EAA/Phononics, 2025.
\textit{Non-resonant acoustic metamaterial approach using plasma actuators.}

% ────────────────────────────────────────────────────────────────
% ENTRIES ADDED FOR Ch.20 (Metalloplasmic Life)
% ────────────────────────────────────────────────────────────────

\bibitem{Artin_Algebra}
M.~Artin,
\emph{Algebra},
2nd ed., Pearson, 2010.

\bibitem{Bartel93}
D.~P.~Bartel and J.~W.~Szostak,
``Isolation of new ribozymes from a large pool of random sequences,''
\emph{Science}, vol.~261, pp.~1411--1418, 1993.

\bibitem{Beinert97}
H.~Beinert, R.~H.~Holm, and E.~M\"unck,
``Iron-sulfur clusters: Nature's modular, multipurpose structures,''
\emph{Science}, vol.~277, no.~5326, pp.~653--659, 1997.

\bibitem{BindiSteinhardt09}
L.~Bindi, P.~J.~Steinhardt, N.~Yao, and P.~J.~Lu,
``Natural quasicrystals,''
\emph{Science}, vol.~324, no.~5932, pp.~1306--1309, 2009.

\bibitem{Birch52}
F.~Birch,
``Elasticity and constitution of the earth's interior,''
\emph{J.\ Geophys.\ Res.}, vol.~57, no.~2, pp.~227--286, 1952.

\bibitem{Breslow59}
R.~Breslow,
``On the mechanism of the formose reaction,''
\emph{Tetrahedron Letters}, vol.~1, no.~21, pp.~22--26, 1959.

\bibitem{Brostigen69}
G.~Brostigen and A.~Kjekshus,
``Redetermined crystal structure of FeS$_2$ (pyrite),''
\emph{Acta Chem.\ Scand.}, vol.~23, pp.~2186--2188, 1969.

\bibitem{Burgess96}
B.~K.~Burgess and D.~J.~Lowe,
``Mechanism of Molybdenum Nitrogenase,''
\emph{Chemical Reviews}, vol.~96, no.~7, pp.~2983--3012, 1996.

\bibitem{Butlerow1861}
A.~Butlerow,
``Bildung einer zuckerartigen Substanz durch Synthese,''
\emph{Justus Liebigs Annalen der Chemie}, vol.~120, pp.~295--298, 1861.
\textit{Discovery of the formose reaction.}

\bibitem{Eck66}
R.~V.~Eck and M.~O.~Dayhoff,
``Evolution of the structure of ferredoxin based on living relics of primitive amino acid sequences,''
\emph{Science}, vol.~152, no.~3720, pp.~363--366, 1966.

\bibitem{Gaucher08}
E.~A.~Gaucher, S.~Govindarajan, and O.~K.~Ganesh,
``Palaeotemperature trend for Precambrian life inferred from resurrected proteins,''
\emph{Nature}, vol.~451, pp.~704--707, 2008.

\bibitem{Gilbert86}
W.~Gilbert,
``Origin of life: The RNA world,''
\emph{Nature}, vol.~319, p.~618, 1986.

\bibitem{Glatzmaier95}
G.~A.~Glatzmaier and P.~H.~Roberts,
``A three-dimensional self-consistent computer simulation of a geomagnetic field reversal,''
\emph{Nature}, vol.~377, pp.~203--209, 1995.

\bibitem{Hall71}
D.~O.~Hall, R.~Cammack, and K.~K.~Rao,
``Role for ferredoxins in the origin of life and biological evolution,''
\emph{Nature}, vol.~233, pp.~136--138, 1971.

\bibitem{Hein00}
J.~E.~Hein, E.~Tse, and D.~G.~Blackmond,
``A route to enantiopure RNA precursors from nearly racemic starting materials,''
\emph{Nature Chemistry}, vol.~3, pp.~704--706, 2011.
\textit{(Key cited as Hein00 in ch.20; year approximate.)}

\bibitem{Hille96}
B.~Hille,
\emph{Ion Channels of Excitable Membranes},
3rd ed., Sinauer Associates, 2001.

\bibitem{Hosoya71}
H.~Hosoya,
``Topological Index: A Newly Proposed Quantity Characterizing the Topological Nature of Structural Isomers of Saturated Hydrocarbons,''
\emph{Bull.\ Chem.\ Soc.\ Japan}, vol.~44, no.~9, pp.~2332--2339, 1971.

\bibitem{Hosoya76}
H.~Hosoya,
``Fibonacci numbers and Fibonacci cubes,''
\emph{Fibonacci Quarterly}, vol.~14, pp.~173--181, 1976.

\bibitem{Huber97}
C.~Huber and G.~W\"achtersh\"auser,
``Activated Acetic Acid by Carbon Fixation on (Fe,Ni)S Under Primordial Conditions,''
\emph{Science}, vol.~276, pp.~245--247, 1997.

\bibitem{Huber98}
C.~Huber and G.~W\"achtersh\"auser,
``Peptides by activation of amino acids with CO on (Ni,Fe)S surfaces:
Implications for the origin of life,''
\emph{Science}, vol.~281, pp.~670--672, 1998.

\bibitem{Johnston01}
D.~T.~Johnston, F.~Wolfe-Simon, A.~Pearson, and A.~H.~Knoll,
``Anoxygenic photosynthesis modulated Proterozoic oxygen and sustained Earth's middle age,''
\emph{Proc.\ Natl.\ Acad.\ Sci.\ USA}, vol.~106, pp.~16925--16929, 2009.
\textit{(Key cited as Johnston01 in ch.20.)}

\bibitem{KKT83}
M.~Kohmoto, L.~P.~Kadanoff, and C.~Tang,
``Localization problem in one dimension: Mapping and escape,''
\emph{Phys.\ Rev.\ Lett.}, vol.~50, no.~23, pp.~1870--1872, 1983.
\textit{Fibonacci chain tight-binding model with Cantor-set energy spectrum.}

\bibitem{Lancaster11}
N.~Lancaster,
``Iron-Sulfur Clusters in Biological Electron Transfer and Sensing,''
in \emph{Comprehensive Inorganic Chemistry II}, Elsevier, 2013.

\bibitem{Lancaster14}
N.~Lancaster et al.,
``X-ray emission spectroscopy evidences a central carbon in the FeMo-cofactor of nitrogenase,''
\emph{Science}, vol.~334, pp.~940--943, 2011.
\textit{(Key cited as Lancaster14 in ch.20.)}

\bibitem{Lane10}
N.~Lane, J.~F.~Allen, and W.~Martin,
``How did LUCA make a living? Chemiosmosis in the origin of life,''
\emph{BioEssays}, vol.~32, pp.~271--280, 2010.

\bibitem{LaRue_PP}
D.~La Rue,
\emph{Principia Psychedelia},
Spectrality Institute, 2026.

\bibitem{Martin03}
W.~Martin and M.~J.~Russell,
``On the origins of cells: a hypothesis for the evolutionary transitions from
abiotic geochemistry to chemoautotrophic prokaryotes, and from prokaryotes to
nucleated cells,''
\emph{Philos.\ Trans.\ R.\ Soc.\ B}, vol.~358, pp.~59--85, 2003.

\bibitem{McCollom99}
T.~M.~McCollom and J.~S.~Seewald,
``Carbon isotope composition of organic compounds produced by abiotic synthesis under hydrothermal conditions,''
\emph{Earth Planet.\ Sci.\ Lett.}, vol.~243, pp.~74--84, 2006.
\textit{(Key cited as McCollom99 in ch.20.)}

\bibitem{McDonough95}
W.~F.~McDonough and S.~-S.~Sun,
``The composition of the Earth,''
\emph{Chem.\ Geology}, vol.~120, pp.~223--253, 1995.

\bibitem{Michibata03}
H.~Michibata, T.~Uyama, and K.~Kanamori,
``The accumulation and reduction of vanadium by ascidians,''
in \emph{Vanadium in Biological Systems}, Springer, pp.~149--169, 2003.

\bibitem{Morowitz92}
H.~J.~Morowitz,
\emph{Beginnings of Cellular Life: Metabolism Recapitulates Biogenesis},
Yale University Press, 1992.

\bibitem{Oro61}
J.~Or\'o,
``Mechanism of synthesis of adenine from hydrogen cyanide under possible primitive earth conditions,''
\emph{Nature}, vol.~191, pp.~1193--1194, 1961.

\bibitem{Pasek05}
M.~A.~Pasek and D.~S.~Lauretta,
``Aqueous corrosion of phosphide minerals from iron meteorites: A highly reactive source of prebiotic phosphorus on the surface of the early Earth,''
\emph{Astrobiology}, vol.~5, no.~4, pp.~515--535, 2005.

\bibitem{Russell97}
M.~J.~Russell and A.~J.~Hall,
``The emergence of life from iron monosulphide bubbles at a submarine hydrothermal
redox and pH front,''
\emph{J.\ Geol.\ Soc.\ Lond.}, vol.~154, pp.~377--402, 1997.

\bibitem{Saenger84}
W.~Saenger,
\emph{Principles of Nucleic Acid Structure},
Springer-Verlag, 1984.

\bibitem{Schrodinger44}
E.~Schr\"odinger,
\emph{What is Life? The Physical Aspect of the Living Cell},
Cambridge University Press, 1944.
\textit{Alias for schrodinger44.}

\bibitem{Spatzal11}
T.~Spatzal, M.~Aksoyoglu, L.~Zhang, S.~L.~A.~Andrade, E.~Schleicher, S.~Weber,
D.~C.~Rees, and O.~Einsle,
``Evidence for interstitial carbon in nitrogenase FeMo cofactor,''
\emph{Science}, vol.~334, pp.~940, 2011.

\bibitem{Tarduno15}
J.~A.~Tarduno et al.,
``A Hadean to Paleoarchean geodynamo recorded by single zircon crystals,''
\emph{Science}, vol.~349, pp.~521--524, 2015.

\bibitem{TsaiGuo00}
A.~P.~Tsai and J.~Q.~Guo,
``Stable binary quasicrystals: The Cd-based family,''
\emph{Phil.\ Mag.\ Lett.}, vol.~81, pp.~L123--L128, 2001.

\bibitem{VineMatthews63}
F.~J.~Vine and D.~H.~Matthews,
``Magnetic anomalies over oceanic ridges,''
\emph{Nature}, vol.~199, pp.~947--949, 1963.

\bibitem{VonDamm90}
K.~L.~Von Damm,
``Seafloor hydrothermal activity: Black smoker chemistry and chimneys,''
\emph{Annual Review of Earth and Planetary Sciences}, vol.~18, pp.~173--204, 1990.

\bibitem{Wachtershauser88}
G.~W\"achtersh\"auser,
``Before enzymes and templates: Theory of surface metabolism,''
\emph{Microbiological Reviews}, vol.~52, pp.~452--484, 1988.

\bibitem{Wachtershauser92}
G.~W\"achtersh\"auser,
``Groundworks for an evolutionary biochemistry: The iron-sulphur world,''
\emph{Prog.\ Biophys.\ Mol.\ Biol.}, vol.~58, pp.~85--201, 1992.

\bibitem{Weiss16}
M.~C.~Weiss, F.~L.~Sousa, N.~Mrnjavac, S.~Neukirchen, M.~Roettger,
S.~Nelson-Sathi, and W.~F.~Martin,
``The physiology and habitat of the last universal common ancestor,''
\emph{Nature Microbiology}, vol.~1, art.~16116, 2016.

% ────────────────────────────────────────────────────────────────
% ENTRIES ADDED FOR Ch.21 (Quasicrystal Physics)
% ────────────────────────────────────────────────────────────────

\bibitem{Dumitrescu22}
P.~Dumitrescu, J.~Bohnet, J.~Gaebler, A.~Hankin, D.~Hayes, A.~Kumar,
B.~Neyenhuis, R.~Vasseur, and A.~Potter,
``Dynamical topological phase realized in a trapped-ion quantum simulator,''
\emph{Nature}, vol.~607, pp.~463--467, 2022.
\textit{Experimental observation of a prethermal discrete time crystal on 57 qubits
using quasi-periodic (Fibonacci) driving.}

\bibitem{Engel07}
G.~S.~Engel, T.~R.~Calhoun, E.~L.~Read, T.-K.~Ahn, T.~Man\v{c}al, Y.-C.~Cheng,
R.~E.~Blankenship, and G.~R.~Fleming,
``Evidence for wavelike energy transfer through quantum coherence in photosynthetic systems,''
\emph{Nature}, vol.~446, pp.~782--786, 2007.
\textit{First evidence of quantum coherence in biological energy transfer (FMO complex).}

\bibitem{HammesSchiffer10}
S.~Hammes-Schiffer and A.~A.~Stuchebrukhov,
``Theory of Coupled Electron and Proton Transfer Reactions,''
\emph{Chemical Reviews}, vol.~110, no.~12, pp.~6939--6960, 2010.

\bibitem{Man05}
W.~Man, M.~Megens, P.~J.~Steinhardt, and P.~M.~Chaikin,
``Experimental measurement of the photonic properties of icosahedral quasicrystals,''
\emph{Nature}, vol.~436, pp.~993--996, 2005.

\bibitem{Marcus56}
R.~A.~Marcus,
``On the Theory of Oxidation-Reduction Reactions Involving Electron Transfer.\ I.,''
\emph{J.\ Chem.\ Phys.}, vol.~24, no.~5, pp.~966--978, 1956.
\textit{Nobel Prize in Chemistry 1992.}

\bibitem{Popp79}
F.-A.~Popp,
"Coherent photon storage of biological systems",
in \emph{Electromagnetic Bio-Information} (F.-A.~Popp, ed.),
Urban \& Schwarzenberg, M\"{u}nchen, 1979, pp.~123-149.
\textit{Foundational work on ultraweak coherent photon emission from living cells.}

\bibitem{Tsai04}
A.~P.~Tsai,
``Discovery of stable icosahedral quasicrystals: progress in understanding structure and properties,''
\emph{Chem.\ Soc.\ Rev.}, vol.~42, pp.~5352--5363, 2013.
\textit{Review of Tsai-type cluster structures in stable icosahedral QCs.}

\bibitem{Vardeny13}
Z.~V.~Vardeny, A.~Nahata, and A.~Agrawal,
``Optics of photonic quasicrystals,''
\emph{Nature Photonics}, vol.~7, pp.~177--187, 2013.

\bibitem{Wilczek12}
F.~Wilczek,
``Quantum Time Crystals,''
\emph{Physical Review Letters}, vol.~109, 160401, 2012.
\textit{Original proposal for discrete time crystals. Alias for wilczek2012.}


% === ADDED: Landmark citations for Force × Matter sections ===

\bibitem{FraustoDaSilva2001}
Frausto da Silva, J.~J.~R. \& Williams, R.~J.~P. (2001).
\textit{The Biological Chemistry of the Elements: The Inorganic Chemistry of Life}.
2nd ed., Oxford University Press.
\textit{The definitive reference for biological essentiality of chemical elements.}

\bibitem{Fukui2004Ar41}
Fukui, M. et al. (2004).
``Twenty years of experience in monitoring ${}^{41}$Ar in a research reactor and decrease of its discharge into the environment.''
\textit{Health Phys.} \textbf{87}(5), 511--519.

\bibitem{Ohring2002ThinFilms}
Ohring, M. (2002).
\textit{Materials Science of Thin Films: Deposition and Structure}.
2nd ed., Academic Press. \S5.7: Sputtering processes.

\bibitem{BurtonKrebs1953}
Burton, K. \& Krebs, H.~A. (1953).
``The free-energy changes associated with the individual steps of the tricarboxylic acid cycle.''
\textit{Biochem. J.} \textbf{54}(1), 94--107.

\end{thebibliography}

\bibitem{univalent_foundations}
The Univalent Foundations Program.
\newblock \emph{Homotopy Type Theory: Univalent Foundations of Mathematics}.
\newblock Institute for Advanced Study, 2013.

\bibitem{tsai_alcufe}
Tsai, A.P., Inoue, A. and Masumoto, T.
\newblock \emph{A Stable Quasicrystal in Al-Cu-Fe System}.
\newblock Japanese Journal of Applied Physics, 1987.

\bibitem{johnson1966}
Johnson, N.W.
\newblock \emph{Convex polyhedra with regular faces}.
\newblock Canadian Journal of Mathematics, 1966.

\bibitem{russell_universal}
Russell, W.
\newblock \emph{The Universal One}.
\newblock Brieger Press, 1926.

\end{document}
